% Options for packages loaded elsewhere
\PassOptionsToPackage{unicode}{hyperref}
\PassOptionsToPackage{hyphens}{url}
%
\documentclass[
  a4paper,
]{article}
\usepackage{amsmath,amssymb}
\usepackage{iftex}
\ifPDFTeX
  \usepackage[T1]{fontenc}
  \usepackage[utf8]{inputenc}
  \usepackage{textcomp} % provide euro and other symbols
\else % if luatex or xetex
  \usepackage{unicode-math} % this also loads fontspec
  \defaultfontfeatures{Scale=MatchLowercase}
  \defaultfontfeatures[\rmfamily]{Ligatures=TeX,Scale=1}
\fi
\usepackage{lmodern}
\ifPDFTeX\else
  % xetex/luatex font selection
  \setmainfont[]{DejaVu Serif}
  \setsansfont[]{DejaVu Sans}
  \setmonofont[]{DejaVu Sans Mono}
\fi
% Use upquote if available, for straight quotes in verbatim environments
\IfFileExists{upquote.sty}{\usepackage{upquote}}{}
\IfFileExists{microtype.sty}{% use microtype if available
  \usepackage[]{microtype}
  \UseMicrotypeSet[protrusion]{basicmath} % disable protrusion for tt fonts
}{}
\makeatletter
\@ifundefined{KOMAClassName}{% if non-KOMA class
  \IfFileExists{parskip.sty}{%
    \usepackage{parskip}
  }{% else
    \setlength{\parindent}{0pt}
    \setlength{\parskip}{6pt plus 2pt minus 1pt}}
}{% if KOMA class
  \KOMAoptions{parskip=half}}
\makeatother
\usepackage{xcolor}
\usepackage[margin=2.5cm]{geometry}
\usepackage{longtable,booktabs,array}
\usepackage{calc} % for calculating minipage widths
% Correct order of tables after \paragraph or \subparagraph
\usepackage{etoolbox}
\makeatletter
\patchcmd\longtable{\par}{\if@noskipsec\mbox{}\fi\par}{}{}
\makeatother
% Allow footnotes in longtable head/foot
\IfFileExists{footnotehyper.sty}{\usepackage{footnotehyper}}{\usepackage{footnote}}
\makesavenoteenv{longtable}
\setlength{\emergencystretch}{3em} % prevent overfull lines
\providecommand{\tightlist}{%
  \setlength{\itemsep}{0pt}\setlength{\parskip}{0pt}}
\setcounter{secnumdepth}{-\maxdimen} % remove section numbering
\ifLuaTeX
  \usepackage{selnolig}  % disable illegal ligatures
\fi
\IfFileExists{bookmark.sty}{\usepackage{bookmark}}{\usepackage{hyperref}}
\IfFileExists{xurl.sty}{\usepackage{xurl}}{} % add URL line breaks if available
\urlstyle{same}
\hypersetup{
  pdfauthor={Matheus Turra Malucelli},
  hidelinks,
  pdfcreator={LaTeX via pandoc}}

\author{Matheus Turra Malucelli}
\date{2026}

\begin{document}

\hypertarget{negociauxe7uxe3o-pruxf3pria-do-emissor-sob-assimetria-informacional}{%
\section{Negociação própria do emissor sob assimetria
informacional}\label{negociauxe7uxe3o-pruxf3pria-do-emissor-sob-assimetria-informacional}}

\hypertarget{uma-arquitetura-regulatuxf3ria-de-capacidade-disclosure-e-aprendizagem-de-mercado}{%
\subsection{Uma arquitetura regulatória de capacidade, disclosure e
aprendizagem de
mercado}\label{uma-arquitetura-regulatuxf3ria-de-capacidade-disclosure-e-aprendizagem-de-mercado}}

\textbf{Autor:} Matheus Turra Malucelli\\
\textbf{Ano:} 2026

\hypertarget{resumo}{%
\subsubsection{Resumo}\label{resumo}}

Esta tese investiga se uma companhia aberta poderia negociar ações de
própria emissão mesmo enquanto possui vantagem informacional sobre si
mesma sem receber capacidade ilimitada e opaca de explorar essa
vantagem. A proposta não busca eliminar a assimetria informacional, mas
disciplinar sua utilização por meio de limites distintos de intervenção
e posição, disclosure condicionado à utilização da capacidade, memória
informacional e aprendizagem do mercado. O desenho persegue
simultaneamente três ganhos institucionais: melhorar a alocação de
capital, tornar observável a expressão econômica da informação privada
do emissor e oferecer maior previsibilidade jurídica para negociação de
boa-fé dentro de um regime especial verificável. A arquitetura é
deliberadamente parametrizável e falsificável: sua contribuição é
decompor o problema em objetos econômicos e regulatórios explícitos,
distinguir invariantes de hipóteses de especificação e tornar possível
testar se existe alguma configuração em que a capacidade economicamente
relevante do emissor seja compatível com integridade, liquidez e
descoberta descentralizada de preço.

\textbf{Palavras-chave:} negociação de ações próprias; recompra de
ações; assimetria informacional; microestrutura de mercado; disclosure;
regulação financeira; descoberta de preço.

© 2026 Matheus Turra Malucelli. Este texto é disponibilizado sob a
licença
\href{https://creativecommons.org/licenses/by/4.0/deed.pt-br}{Creative
Commons Atribuição 4.0 Internacional --- CC BY 4.0}. A licença exige
atribuição adequada, indicação de alterações e referência à licença.

\begin{center}\rule{0.5\linewidth}{0.5pt}\end{center}

\hypertarget{sumuxe1rio}{%
\subsection{Sumário}\label{sumuxe1rio}}

\hypertarget{i-problema-tese-e-escopo}{%
\subsubsection{\texorpdfstring{\protect\hyperlink{i-problema-tese-e-escopo}{I.
Problema, tese e
escopo}}{I. Problema, tese e escopo}}\label{i-problema-tese-e-escopo}}

\begin{enumerate}
\def\labelenumi{\arabic{enumi}.}
\tightlist
\item
  \protect\hyperlink{1-a-provocacao-institucional}{A provocação
  institucional}
\item
  \protect\hyperlink{2-a-tese-economica}{A tese econômica}
\item
  \protect\hyperlink{3-ganhos-pretendidos-alocacao-informacao-e-seguranca-juridica}{Ganhos
  pretendidos: alocação, informação e segurança jurídica}
\item
  \protect\hyperlink{4-o-que-o-protocolo-nao-pretende-fazer}{O que o
  protocolo não pretende fazer}
\end{enumerate}

\hypertarget{ii-ontologia-econuxf4mica-do-mecanismo}{%
\subsubsection{\texorpdfstring{\protect\hyperlink{ii-ontologia-economica-do-mecanismo}{II.
Ontologia econômica do
mecanismo}}{II. Ontologia econômica do mecanismo}}\label{ii-ontologia-econuxf4mica-do-mecanismo}}

\begin{enumerate}
\def\labelenumi{\arabic{enumi}.}
\setcounter{enumi}{4}
\tightlist
\item
  \protect\hyperlink{5-a-unidade-regulada-e-o-estado-do-emissor}{A
  unidade regulada e o estado do emissor}
\item
  \protect\hyperlink{6-gross-e-net-intervencao-e-posicao}{Gross e Net:
  intervenção e posição}
\item
  \protect\hyperlink{7-treasury-e-o-lado-sell}{Treasury e o lado SELL}
\item
  \protect\hyperlink{8-h-capacidade-especial-de-negociacao}{H:
  capacidade especial de negociação}
\item
  \protect\hyperlink{9-h0-hmax-e-a-escala-da-capacidade}{H₀,
  \(H_{\max}\) e a escala da capacidade}
\item
  \protect\hyperlink{10-q-quando-a-utilizacao-da-capacidade-obriga-disclosure}{Q:
  alerta quantitativo e avaliação de disclosure}
\item
  \protect\hyperlink{11-u-utilizacao-como-intensidade-observavel-do-sinal}{U:
  utilização como intensidade observável do sinal}
\item
  \protect\hyperlink{12-informacao-produzida-pelo-emissor-informacao-do-mercado-e-atividade-gross}{Informação
  produzida pelo emissor, informação do mercado e atividade Gross}
\item
  \protect\hyperlink{13-decay-memoria-e-saturacao}{Decay, memória e
  saturação}
\end{enumerate}

\hypertarget{iii-tempo-disclosure-e-transiuxe7uxe3o-de-estado}{%
\subsubsection{\texorpdfstring{\protect\hyperlink{iii-tempo-disclosure-e-transicao-de-estado}{III.
Tempo, disclosure e transição de
estado}}{III. Tempo, disclosure e transição de estado}}\label{iii-tempo-disclosure-e-transiuxe7uxe3o-de-estado}}

\begin{enumerate}
\def\labelenumi{\arabic{enumi}.}
\setcounter{enumi}{13}
\tightlist
\item
  \protect\hyperlink{14-w-janela-regulatoria-disclosure-e-transicao-de-estado}{W:
  janela regulatória, disclosure e transição de estado}
\item
  \protect\hyperlink{15-modelo-temporal-rolling}{Modelo temporal
  rolling}
\item
  \protect\hyperlink{16-modelo-temporal-calendarizado}{Modelo temporal
  calendarizado}
\item
  \protect\hyperlink{17-o-que-realmente-diferencia-as-duas-branches}{Duas
  topologias temporais para o mesmo mecanismo}
\item
  \protect\hyperlink{18-cooldown-reversao-e-fato-relevante}{Intervalos
  de reação, reversão e Fato Relevante}
\end{enumerate}

\hypertarget{iv-informauxe7uxe3o-aprendizagem-e-integridade}{%
\subsubsection{\texorpdfstring{\protect\hyperlink{iv-informacao-aprendizagem-e-integridade}{IV.
Informação, aprendizagem e
integridade}}{IV. Informação, aprendizagem e integridade}}\label{iv-informauxe7uxe3o-aprendizagem-e-integridade}}

\begin{enumerate}
\def\labelenumi{\arabic{enumi}.}
\setcounter{enumi}{18}
\tightlist
\item
  \protect\hyperlink{19-a-companhia-como-trader-identificavel-e-potencialmente-informado}{A
  companhia como trader identificável e potencialmente informado}
\item
  \protect\hyperlink{20-reputacao-antecipacao-e-reflexividade}{Reputação,
  antecipação e reflexividade}
\item
  \protect\hyperlink{21-o-risco-de-o-sinal-quebrar-o-mercado}{Saliência,
  familiaridade e aprendizagem do sinal}
\item
  \protect\hyperlink{22-manipulacao-gross-e-transferencia-de-valor-a-terceiros}{Captura
  parcial de valor, fabricação de sinal e fronteira da manipulação}
\item
  \protect\hyperlink{23-disciplina-endogena-e-enforcement-ex-post}{Disciplina
  endógena, desenho adversarial e enforcement ex post}
\end{enumerate}

\hypertarget{v-especificauxe7uxe3o-calibrauxe7uxe3o-e-laboratuxf3rio-regulatuxf3rio}{%
\subsubsection{\texorpdfstring{\protect\hyperlink{v-especificacao-calibracao-e-laboratorio-regulatorio}{V.
Especificação, calibração e laboratório
regulatório}}{V. Especificação, calibração e laboratório regulatório}}\label{v-especificauxe7uxe3o-calibrauxe7uxe3o-e-laboratuxf3rio-regulatuxf3rio}}

\begin{enumerate}
\def\labelenumi{\arabic{enumi}.}
\setcounter{enumi}{23}
\tightlist
\item
  \protect\hyperlink{24-arquitetura-especificacao-e-calibracao}{Arquitetura,
  especificação e calibração}
\item
  \protect\hyperlink{25-observabilidade-externa-e-flexibilidade-operacional}{Observabilidade
  externa e flexibilidade operacional}
\item
  \protect\hyperlink{26-variaveis-modulares-e-variable-builder}{Variáveis
  modulares e Variable Builder}
\item
  \protect\hyperlink{27-invariantes-e-espaco-experimental}{Invariantes e
  espaço experimental}
\item
  \protect\hyperlink{28-criterio-de-viabilidade}{Critério de
  viabilidade}
\end{enumerate}

\hypertarget{vi-eviduxeancia-precedentes-e-validauxe7uxe3o}{%
\subsubsection{\texorpdfstring{\protect\hyperlink{vi-evidencia-precedentes-e-validacao}{VI.
Evidência, precedentes e
validação}}{VI. Evidência, precedentes e validação}}\label{vi-eviduxeancia-precedentes-e-validauxe7uxe3o}}

\begin{enumerate}
\def\labelenumi{\arabic{enumi}.}
\setcounter{enumi}{28}
\tightlist
\item
  \protect\hyperlink{29-o-que-a-literatura-ja-permite-afirmar}{O que a
  literatura já permite afirmar}
\item
  \protect\hyperlink{30-o-que-precedentes-regulatorios-sustentam-e-o-que-nao-sustentam}{O
  que precedentes regulatórios sustentam --- e o que não sustentam}
\item
  \protect\hyperlink{31-casos-administrativos-como-testes-de-realidade}{Casos
  administrativos como testes de realidade}
\item
  \protect\hyperlink{32-estrategia-de-validacao-e-laboratorio-multiagente}{Estratégia
  de validação e laboratório multiagente}
\end{enumerate}

\hypertarget{vii-conclusuxe3o}{%
\subsubsection{\texorpdfstring{\protect\hyperlink{vii-conclusao}{VII.
Conclusão}}{VII. Conclusão}}\label{vii-conclusuxe3o}}

\begin{enumerate}
\def\labelenumi{\arabic{enumi}.}
\setcounter{enumi}{32}
\tightlist
\item
  \protect\hyperlink{33-o-que-esta-tese-estabelece}{O que esta tese
  estabelece}
\item
  \protect\hyperlink{34-conclusao-conceitual}{Conclusão conceitual}
\item
  \protect\hyperlink{35-referencias}{Referências}
\end{enumerate}

\begin{center}\rule{0.5\linewidth}{0.5pt}\end{center}

\hypertarget{i-problema-tese-e-escopo}{}

\hypertarget{i-problema-tese-e-escopo-1}{%
\section{I. Problema, tese e escopo}\label{i-problema-tese-e-escopo-1}}

\hypertarget{1-a-provocacao-institucional}{}

\hypertarget{1-a-provocauxe7uxe3o-institucional}{%
\subsection{1. A provocação
institucional}\label{1-a-provocauxe7uxe3o-institucional}}

Uma companhia aberta pode conhecer mudanças relevantes em seu próprio
estado econômico antes que o restante do mercado consiga incorporá-las
ao preço. Essa assimetria não é uma anomalia: ela decorre do fato de que
a companhia produz, observa e interpreta informação interna sobre si
mesma.

O problema aparece quando essa companhia quer negociar ações de própria
emissão. Imagine que WXYZ3 esteja sendo negociada a R\$ 20 e que a
administração, por razões que ainda não são públicas, decida comprar. No
mercado ordinário, os demais participantes observam fluxo comprador, mas
não necessariamente sabem que a própria WXYZ está por trás dele.

A intervenção institucional proposta é simples em sua origem:
\textbf{quando o emissor negociar a própria ação dentro do regime
especial, sua atuação deixa de ser anônima}. O mercado secundário
continua sendo o mesmo; não se cria um book separado. O que muda é que a
participação do emissor passa a ser atribuível e submetida a limites
próprios.

Em termos intuitivos, o mercado poderia saber que:

\begin{quote}
\textbf{COMPANHIA WXYZ --- BUY --- quantidade executada --- preço}
\end{quote}

ou que:

\begin{quote}
\textbf{COMPANHIA WXYZ --- SELL --- quantidade executada --- preço}
\end{quote}

Isso não revela a informação privada que motivou a decisão. Revela outra
coisa: \textbf{a própria companhia decidiu comprometer capital naquela
direção}. O mercado continua livre para interpretar esse comportamento
como sinal informacional, alocação de caixa, necessidade de
financiamento, erro de administração ou qualquer outra hipótese
compatível com os dados.

Essa camada não substitui Fato Relevante. Obrigações de divulgação
material continuam existindo segundo sua própria lógica. A operação
identificada responde a uma pergunta diferente: não necessariamente
\emph{o que aconteceu dentro da companhia?}, mas \emph{o que a companhia
decidiu fazer economicamente?}

O regime jurídico brasileiro trata negociação de ações de própria
emissão e negociação na pendência de informação relevante não divulgada
como matérias reguladas, hoje principalmente pelas Resoluções CVM 77 e
44 (\protect\hyperlink{ref-1}{CVM, 2022}; \protect\hyperlink{ref-2}{CVM,
2021}). Esta tese não descreve o regime vigente como autorização para a
companhia ``fazer insider trading''. A proposta é deliberadamente
contrafactual: pergunta como deveria ser desenhada uma exceção
regulatória específica se quiséssemos permitir que o emissor utilizasse
economicamente parte de sua vantagem informacional sem conceder a ele
capacidade ilimitada e opaca de exploração dessa vantagem.

A pergunta deixa de ser apenas:

\begin{quote}
\textbf{Quando o emissor deve ser impedido de negociar?}
\end{quote}

para se tornar:

\begin{quote}
\textbf{Existe uma arquitetura em que o emissor possa negociar mesmo
possuindo vantagem informacional sobre si próprio, preservando
capacidade economicamente relevante de agir sem excluir indefinidamente
o restante do mercado da descoberta de preço?}
\end{quote}

A mudança altera o objeto regulatório. Em vez de tentar eliminar a
assimetria antes da negociação, o protocolo tenta controlar \textbf{como
ela pode ser exercida, quanto pode ser exercida e que informação a
própria utilização dessa capacidade deve devolver ao mercado}.

A cadeia causal que será formalizada nas seções seguintes é:

\begin{quote}
convicção privada → decisão econômica → execução identificável →
utilização mensurável de capacidade → eventual disclosure formal →
atualização do estado informacional → reação do mercado
\end{quote}

A companhia continua sabendo mais. A diferença é que sua ação deixa de
ser silenciosa e ilimitada.

Essa é a premissa institucional que orienta o restante da tese: a
companhia pode ser a primeira a saber e a primeira a agir, mas não deve
conseguir capturar sozinha todo o ajuste de preço produzido por sua
vantagem. Sua execução pode iniciar a descoberta; a escala, o disclosure
e o tempo precisam preservar espaço para que o mercado também reaja,
interprete e negocie.

\begin{center}\rule{0.5\linewidth}{0.5pt}\end{center}

\hypertarget{2-a-tese-economica}{}

\hypertarget{2-a-tese-econuxf4mica}{%
\subsection{2. A tese econômica}\label{2-a-tese-econuxf4mica}}

A hipótese central pode ser expressa sem álgebra sofisticada:

\begin{quote}
\textbf{A vantagem informacional não é, por si só, o objeto que o
mecanismo tenta eliminar. O objeto é a exclusividade com que essa
vantagem pode ser monetizada.}
\end{quote}

A companhia pode ser a primeira a saber e a primeira a agir, mas não
deve receber capacidade indefinida para ser a única participante da
descoberta de preço que sua própria informação produz.

A arquitetura procura, portanto, combinar três propriedades:

\begin{enumerate}
\def\labelenumi{\arabic{enumi}.}
\tightlist
\item
  capacidade economicamente relevante para o emissor;
\item
  produção progressiva de informação observável para o mercado;
\item
  segurança jurídica ex ante para operações de boa-fé, condicionada a
  limites verificáveis de intervenção, posição e disclosure.
\end{enumerate}

A existência simultânea dessas propriedades é hipótese, não conclusão.

Chamemos de \textbf{Θ} o conjunto de configurações possíveis do
mecanismo e de \textbf{Θᵥ} a região que satisfaz critérios mínimos de
viabilidade econômica e integridade de mercado. A pergunta é
simplesmente:

\begin{quote}
\textbf{Θᵥ é vazio ou existe pelo menos uma configuração viável?}
\end{quote}

Nada nesta tese exige que Θᵥ seja diferente de vazio. Se toda
configuração que preserve utilidade econômica para a companhia produzir
deterioração inaceitável de liquidez, manipulação, reflexividade ou
transferência de valor, a hipótese institucional deve ser rejeitada.

\begin{center}\rule{0.5\linewidth}{0.5pt}\end{center}

\hypertarget{3-ganhos-pretendidos-alocacao-informacao-e-seguranca-juridica}{}

\hypertarget{3-ganhos-pretendidos-alocauxe7uxe3o-informauxe7uxe3o-e-seguranuxe7a-juruxeddica}{%
\subsection{3. Ganhos pretendidos: alocação, informação e segurança
jurídica}\label{3-ganhos-pretendidos-alocauxe7uxe3o-informauxe7uxe3o-e-seguranuxe7a-juruxeddica}}

A proposta só é economicamente interessante se entregar algo além de
transparência formal. Ela busca combinar três ganhos que hoje entram em
tensão: \textbf{capacidade de alocação de capital pelo emissor, produção
de informação para o mercado e maior previsibilidade jurídica para a
própria negociação}.

O primeiro ganho é permitir que a companhia coloque patrimônio
corporativo em uma decisão sobre o próprio ativo. Se comprar barato, o
benefício pertence à companhia e, indiretamente, aos seus acionistas; se
comprar caro, o patrimônio corporativo absorve a perda. No lado
vendedor, a companhia pode reduzir Treasury quando essa decisão fizer
sentido econômico, sempre limitada ao estoque de ações que efetivamente
possui. A tese não presume que BUY signifique subavaliação nem que SELL
signifique sobreavaliação. Essas interpretações pertencem ao mercado.

O segundo ganho é informacional. A operação deixa um rastro atribuível:
direção, quantidade, preço, utilização da capacidade e histórico. O
emissor não precisa revelar seu raciocínio privado para que sua decisão
econômica se torne um dado adicional da descoberta de preço. A proposta
procura fazer com que o mercado não seja completamente excluído enquanto
a companhia exerce sua vantagem.

O terceiro ganho é jurídico. A segurança jurídica é um pilar, não um
efeito colateral. O regime especial pretende substituir parte de uma
fronteira dependente de estado informacional e interpretação ex post por
uma fronteira operacional mais explícita: capacidade disponível,
utilização, triggers de disclosure, posição de Treasury e invariantes
verificáveis. A companhia que opere de boa-fé dentro do regime deve
conseguir saber ex ante quais condições tornam sua atuação admissível.

Isso não cria imunidade. Fraude, manipulação, coordenação artificial,
informação falsa e outras condutas abusivas permanecem sujeitas a
supervisão e julgamento. O safe harbour proposto é \textbf{condicionado
ao protocolo}: reduz incerteza sobre a possibilidade de o emissor
negociar enquanto possui vantagem informacional, sem transformar
conformidade mecânica em licença para manipular o mercado.

Esses três ganhos podem entrar em conflito. Capacidade demais pode
excluir o mercado; disclosure demais ou cedo demais pode tornar a
capacidade economicamente inútil; limites mal desenhados podem criar
estratégias mecânicas de exploração; segurança jurídica excessivamente
ampla pode proteger condutas que deveriam continuar investigáveis. A
arquitetura existe para tornar esses trade-offs explícitos e testáveis.

\begin{center}\rule{0.5\linewidth}{0.5pt}\end{center}

\hypertarget{4-o-que-o-protocolo-nao-pretende-fazer}{}

\hypertarget{4-o-que-o-protocolo-nuxe3o-pretende-fazer}{%
\subsection{4. O que o protocolo não pretende
fazer}\label{4-o-que-o-protocolo-nuxe3o-pretende-fazer}}

O protocolo não pretende tornar participantes igualmente informados,
igualmente inteligentes ou igualmente capazes de interpretar fluxo.

Ele também não pretende substituir as regras de Fato Relevante,
legalizar negociação de administradores ou controladores pessoas
físicas, eliminar investigação de manipulação ou declarar legítima
qualquer operação realizada pela companhia.

O regime especial seria restrito ao emissor negociando ações de própria
emissão dentro de regras públicas específicas.

A companhia não seria obrigada a divulgar imediatamente \textbf{por que}
decidiu comprar ou vender. Seu raciocínio permanece privado. O mercado
recebe sua ação, a intensidade dessa ação e o histórico produzido pelo
protocolo.

A participação no regime-base pressupõe ainda um perímetro instrumental
estreito. A companhia não poderia combinar a capacidade especial sobre a
ação com exposição derivativa ao preço da própria ação. A razão é
econômica: permitir simultaneamente capacidade regulatória especial
sobre o subjacente e payoff alavancado ou não linear sobre o próprio
sinal abriria um segundo canal de monetização não disciplinado pela
posição física.

Essa proibição é uma escolha conservadora de desenho. O safe harbour
europeu para programas de recompra também delimita sua exceção à
negociação das próprias ações sob condições específicas; isso serve como
precedente de perímetro estreito, não como demonstração de que nossa
proibição seja ótima (\protect\hyperlink{ref-3}{European Commission,
2016}).

\begin{center}\rule{0.5\linewidth}{0.5pt}\end{center}

\hypertarget{ii-ontologia-economica-do-mecanismo}{}

\hypertarget{ii-ontologia-econuxf4mica-do-mecanismo-1}{%
\section{II. Ontologia econômica do
mecanismo}\label{ii-ontologia-econuxf4mica-do-mecanismo-1}}

\hypertarget{5-a-unidade-regulada-e-o-estado-do-emissor}{}

\hypertarget{5-a-unidade-regulada-e-o-estado-do-emissor}{%
\subsection{5. A unidade regulada e o estado do
emissor}\label{5-a-unidade-regulada-e-o-estado-do-emissor}}

A capacidade pertence ao \textbf{ticker/emissor}, não ao setor.

Para cada emissor i e cada janela regulatória n, o protocolo calcula
três capacidades:

\begin{itemize}
\tightlist
\item
  \textbf{\(H_{\mathrm{gross}}\)}: capacidade total de intervenção;
\item
  \textbf{\(H_{\mathrm{net}+}\)}: capacidade de aumentar a posição
  própria;
\item
  \textbf{\(H_{\mathrm{net}-}\)}: capacidade de reduzir a posição
  própria.
\end{itemize}

Essas capacidades são funções de um estado observável e regulatoriamente
definido. Esse estado pode conter características do ativo, posição de
Treasury, memória informacional, utilização histórica e variáveis
externas selecionadas.

A tese não fixa uma única equação para H. Ela fixa o \textbf{significado
econômico} de cada dimensão e as propriedades que qualquer função
escolhida deve preservar.

Essa distinção é importante. O protocolo não precisa decidir hoje quanto
peso atribuir à volatilidade ou à profundidade, nem se determinada
variável externa deve sequer integrar o cálculo. Essas são decisões de
especificação e calibração. A ontologia vem antes.

\begin{center}\rule{0.5\linewidth}{0.5pt}\end{center}

\hypertarget{6-gross-e-net-intervencao-e-posicao}{}

\hypertarget{6-gross-e-net-intervenuxe7uxe3o-e-posiuxe7uxe3o}{%
\subsection{6. Gross e Net: intervenção e
posição}\label{6-gross-e-net-intervenuxe7uxe3o-e-posiuxe7uxe3o}}

Definimos, em cada janela:

\begin{quote}
\textbf{Gross = BUY + SELL}
\end{quote}

\begin{quote}
\textbf{Net = BUY − SELL}
\end{quote}

Exemplo simples:

\begin{longtable}[]{@{}
  >{\raggedleft\arraybackslash}p{(\columnwidth - 6\tabcolsep) * \real{0.2500}}
  >{\raggedleft\arraybackslash}p{(\columnwidth - 6\tabcolsep) * \real{0.2500}}
  >{\raggedleft\arraybackslash}p{(\columnwidth - 6\tabcolsep) * \real{0.2500}}
  >{\raggedleft\arraybackslash}p{(\columnwidth - 6\tabcolsep) * \real{0.2500}}@{}}
\toprule\noalign{}
\begin{minipage}[b]{\linewidth}\raggedleft
BUY
\end{minipage} & \begin{minipage}[b]{\linewidth}\raggedleft
SELL
\end{minipage} & \begin{minipage}[b]{\linewidth}\raggedleft
Gross
\end{minipage} & \begin{minipage}[b]{\linewidth}\raggedleft
Net
\end{minipage} \\
\midrule\noalign{}
\endhead
\bottomrule\noalign{}
\endlastfoot
100 & 0 & 100 & +100 \\
0 & 100 & 100 & −100 \\
100 & 100 & 200 & 0 \\
\end{longtable}

O terceiro caso é o motivo de a separação ser necessária. Net igual a
zero não significa ausência de atuação no mercado.

\textbf{Gross mede intensidade de intervenção.} Uma companhia que compra
e vende repetidamente pode afetar liquidez, fluxo, expectativa e preço
mesmo terminando sem alteração líquida de posição.

\textbf{Net mede alteração de exposição econômica.} É a dimensão mais
diretamente relacionada à captura de vantagem informacional por
construção ou redução de posição.

A associação principal é, portanto:

\begin{quote}
\(H_{\mathrm{gross}}\) → disciplina principalmente intervenção e risco
microestrutural
\end{quote}

\begin{quote}
\(H_{\mathrm{net}}\) → disciplina principalmente alteração posicional
sob vantagem informacional
\end{quote}

As dimensões se sobrepõem, mas não são substitutas.

Uma reversão não restitui Gross. Se a companhia compra 100 e depois
vende 100, ela pode voltar ao Net inicial, mas consumiu 200 unidades de
intervenção. Essa propriedade é uma defesa natural contra estratégias de
alternância: trocar de direção pode reabrir espaço posicional, mas nunca
apaga o custo regulatório da atividade já realizada.

\begin{center}\rule{0.5\linewidth}{0.5pt}\end{center}

\hypertarget{7-treasury-e-o-lado-sell}{}

\hypertarget{7-treasury-e-o-lado-sell}{%
\subsection{7. Treasury e o lado SELL}\label{7-treasury-e-o-lado-sell}}

O regime-base não permite short verdadeiro da companhia em sua própria
ação.

Portanto:

\begin{quote}
\textbf{Treasury ≥ 0}
\end{quote}

\begin{quote}
\textbf{SELL ≤ ações disponíveis em Treasury}
\end{quote}

O lado vendedor significa redução de ações previamente mantidas em
tesouraria, não criação de exposição econômica negativa.

Três grandezas patrimoniais devem permanecer distintas. \textbf{Ações
emitidas} abrangem as ações em circulação e as ações mantidas em
tesouraria. \textbf{Ações em circulação} excluem Treasury. \textbf{Free
float} corresponde apenas à parcela das ações em circulação efetivamente
disponível para negociação dispersa no mercado. Portanto:

\begin{quote}
\textbf{ações emitidas = ações em circulação + Treasury}
\end{quote}

\begin{quote}
\textbf{free float ≤ ações em circulação}
\end{quote}

Em geral, free float somado a Treasury não corresponde ao total emitido,
porque ações em circulação podem permanecer com controladores,
administradores ou outros blocos não integrantes do free float. O
protocolo deve receber essas grandezas como dados separados, com
proveniência e data de referência próprias.

Isso torna \(H_{\mathrm{net}}\) assimétrico por construção. Podemos ter
\(H_{\mathrm{net}+}\) e \(H_{\mathrm{net}-}\) diferentes, porque comprar
e vender partem de restrições físicas distintas.

É útil normalizar a posição de Treasury entre 0 e 1. No extremo
inferior, não existe ação disponível para venda; no extremo superior,
não existe espaço físico adicional para aquisição dentro do limite
societário considerado. Assim:

\begin{quote}
Treasury = 0 → capacidade SELL líquida deve convergir a zero
\end{quote}

\begin{quote}
Treasury = limite máximo → capacidade BUY líquida deve convergir a zero
\end{quote}

A posição física é diferente da informação. Mesmo um sinal informacional
extremamente forte não pode criar ações em Treasury que não existem nem
ampliar um limite societário por construção.

Também não se presume que SELL signifique necessariamente visão negativa
sobre o valor da ação. A companhia pode vender para financiar projetos,
restaurar free float, desfazer recompra anterior ou realocar capital. O
mercado aprende o significado empírico desse comportamento ao longo do
tempo.

\begin{center}\rule{0.5\linewidth}{0.5pt}\end{center}

\hypertarget{8-h-capacidade-especial-de-negociacao}{}

\hypertarget{8-h-capacidade-especial-de-negociauxe7uxe3o}{%
\subsection{8. H: capacidade especial de
negociação}\label{8-h-capacidade-especial-de-negociauxe7uxe3o}}

H é o orçamento regulatório de utilização da vantagem especial concedida
ao emissor.

A arquitetura possui três H porque existem três perguntas diferentes:

\begin{quote}
\textbf{Quanto a companhia pode mexer no próprio mercado?} →
\(H_{\mathrm{gross}}\)
\end{quote}

\begin{quote}
\textbf{Quanto pode aumentar sua posição?} → \(H_{\mathrm{net}+}\)
\end{quote}

\begin{quote}
\textbf{Quanto pode reduzir sua posição?} → \(H_{\mathrm{net}-}\)
\end{quote}

H é um \textbf{hard cap}. A execução não pode ultrapassá-lo dentro da
janela aplicável. Se uma ordem excede a capacidade residual, executa-se
apenas a maior parcela simultaneamente compatível com
\(H_{\mathrm{gross}}\), \(H_{\mathrm{net}+}\), \(H_{\mathrm{net}-}\) e
as restrições físicas aplicáveis. O residual não executado permanece
registrado como ordem rejeitada, acompanhado da quantidade e do limite
vinculante. Solicitação, execução e rejeição são objetos distintos.

A função que calcula H pode ser representada abstratamente como:

\begin{quote}
\textbf{H = função(estado do ticker, posição, informação, variáveis
regulatórias selecionadas)}
\end{quote}

Características setoriais podem integrar os inputs da função de H quando
forem economicamente relevantes. Ainda assim, elas apenas ajudam a
caracterizar o estado do ticker: a capacidade resultante é
individualizada e pertence ao emissor avaliado, não ao setor como
unidade regulada.

A função também não precisa ser puramente multiplicativa. O laboratório
poderá testar multiplicação, adição, funções limitadas, min/max, regras
por faixas ou outras composições, desde que preservem os invariantes da
arquitetura.

Um \textbf{plano público voluntário de recompra} pode ser tratado como
variável candidata de H, e não como obrigação estrutural do regime. A
lógica é limitada: ao anunciar previamente uma intenção agregada --- por
exemplo, adquirir até determinado montante ao longo de um ano --- a
companhia reduz parte da surpresa informacional da primeira execução.
Uma especificação pode reconhecer esse pré-disclosure com aumento
moderado da capacidade inicial. O valor anunciado, porém, \textbf{não se
converte em capacidade automática}. A continuidade da escala depende de
execução efetiva e das demais regras do protocolo. Se o plano aumenta H,
quanto aumenta e sob quais limites são questões de especificação e
calibração.

\begin{center}\rule{0.5\linewidth}{0.5pt}\end{center}

\hypertarget{9-h0-hmax-e-a-escala-da-capacidade}{}

\hypertarget{9-hux2080-h_max-e-a-escala-da-capacidade}{%
\subsection{\texorpdfstring{9. H₀, \(H_{\max}\) e a escala da
capacidade}{9. H₀, H\_\{\textbackslash max\} e a escala da capacidade}}\label{9-hux2080-h_max-e-a-escala-da-capacidade}}

Para tornar a dinâmica inteligível, é útil separar três ideias.

\textbf{H₀} é uma referência estrutural ou neutra para o ticker. Ela
pode depender de características como volume, profundidade, spread, free
float, turnover e volatilidade. H₀ não é um piso garantido.

\textbf{\(H_{\max}\)} é uma capacidade máxima admissível sob determinada
especificação, mantidas constantes as demais restrições. É o teto para o
qual o componente informacional poderia empurrar H --- não uma
capacidade automaticamente concedida.

\textbf{H calculado} é a capacidade produzida pela função para o estado
atual, antes das restrições físicas externas à função.

Uma forma conceitual simples é:

\begin{quote}
\textbf{H calculado = H₀ + escala informacional × (\(H_{\max}\) − H₀)}
\end{quote}

A ``escala informacional'' varia entre uma condição neutra e uma
condição máxima. A função exata não é fixada aqui.

A representação separa duas decisões. Primeiro define-se o intervalo
economicamente admissível entre uma capacidade neutra e uma capacidade
máxima. Depois define-se como o estado informacional move H dentro desse
intervalo.

A capacidade efetivamente disponível resulta da aplicação das restrições
físicas ao H calculado. No lado vendedor:

\begin{quote}
\textbf{\(H_{\mathrm{net}-}\) efetivo = min(\(H_{\mathrm{net}-}\)
calculado, Treasury disponível)}
\end{quote}

No lado comprador, uma especificação pode ainda limitar
\(H_{\mathrm{net}+}\) pelo espaço remanescente até um limite societário
ou regulatório de posição própria. Essa trava adicional pertence ao
espaço experimental; a impossibilidade de vender mais ações do que a
Treasury disponível pertence à arquitetura.

Uma função saturante é uma candidata natural porque impede que H cresça
indefinidamente. Ela também pode representar retornos marginais
decrescentes de sinais repetidos na mesma direção --- vários disclosures
semelhantes podem acrescentar progressivamente menos informação nova
---, mas essa interpretação \textbf{não define} o conceito de saturação.
É uma hipótese adicional sobre a forma do mapeamento entre informação e
capacidade.

H₀ e \(H_{\max}\) permanecem objetos de calibração. Um H₀ muito pequeno
pode tornar o regime pouco útil no curto prazo, embora a necessidade de
entrar gradualmente seja parte deliberada da arquitetura. Essa
progressão não deve responder a um sinal mínimo ou meramente mecânico: a
escala futura deve depender de utilização relevante da capacidade e de
execução observável. Assim, o risco de ``aquecer'' H por pequenas
operações artificiais é tratado pela própria regra de progressão, e não
pela concessão de uma capacidade inicial maior. Um H₀ excessivo, por
outro lado, pode conceder grande vantagem antes que o emissor tenha
produzido informação compensatória. O mesmo vale para a distância entre
H₀ e \(H_{\max}\).

A velocidade de aproximação entre H₀ e \(H_{\max}\) também é
substantiva. Uma progressão lenta demais pode impedir uma companhia que
pretende executar uma alocação economicamente relevante ao longo de
meses de atingir escala útil; uma progressão rápida demais pode permitir
que pequenas operações iniciais fabriquem capacidade desproporcional.
Por isso o laboratório deve testar crescimento linear, acelerado e
saturante, inclusive especificações em que \textbf{persistência
direcional coerente acelera a liberação de capacidade} sem tornar
execução passada um direito automático a H futuro.

Reversões podem receber tratamento assimétrico. É plausível testar se
uma mudança de BUY para SELL degrada parcialmente a capacidade acumulada
na direção anterior e/ou faz \(H_{\mathrm{net}-}\) recomeçar de uma base
menor, preservada a possibilidade legítima de realizar Treasury. Essa
fricção pode dificultar a sequência BUY → reação do mercado → SELL em
escala simétrica imediata. Ela não é canonizada como regra: se
excessiva, pode permitir que o próprio mercado encurrale a companhia ao
explorar sua dificuldade regulatória de saída.

\begin{center}\rule{0.5\linewidth}{0.5pt}\end{center}

\hypertarget{10-q-quando-a-utilizacao-da-capacidade-obriga-disclosure}{}

\hypertarget{10-q-alerta-quantitativo-e-avaliauxe7uxe3o-de-disclosure}{%
\subsection{10. Q: alerta quantitativo e avaliação de
disclosure}\label{10-q-alerta-quantitativo-e-avaliauxe7uxe3o-de-disclosure}}

H responde ``quanto pode executar''. Q responde ``a partir de que
intensidade a atividade deve ser sinalizada ao emissor e submetida à
avaliação quantitativa de disclosure no corte aplicável''.

Para cada dimensão pode existir um par Q/H:

\begin{longtable}[]{@{}
  >{\raggedright\arraybackslash}p{(\columnwidth - 4\tabcolsep) * \real{0.3400}}
  >{\raggedright\arraybackslash}p{(\columnwidth - 4\tabcolsep) * \real{0.3300}}
  >{\raggedright\arraybackslash}p{(\columnwidth - 4\tabcolsep) * \real{0.3300}}@{}}
\toprule\noalign{}
\begin{minipage}[b]{\linewidth}\raggedright
Dimensão
\end{minipage} & \begin{minipage}[b]{\linewidth}\raggedright
Limiar de disclosure
\end{minipage} & \begin{minipage}[b]{\linewidth}\raggedright
Limite rígido
\end{minipage} \\
\midrule\noalign{}
\endhead
\bottomrule\noalign{}
\endlastfoot
intervenção & \(Q_{\mathrm{gross}}\) & \(H_{\mathrm{gross}}\) \\
posição BUY & \(Q_{\mathrm{net}+}\) & \(H_{\mathrm{net}+}\) \\
posição SELL & \(Q_{\mathrm{net}-}\) & \(H_{\mathrm{net}-}\) \\
\end{longtable}

Economicamente, há duas famílias de fronteira --- intervenção e posição
--- e, operacionalmente, a fronteira posicional pode ser direcional.

Q pode ser representado como uma proporção de H:

\begin{quote}
\textbf{Q = ρ × H}, com 0 \textless{} ρ ≤ 1.
\end{quote}

Se ρ = 1, o limiar quantitativo coincide com o esgotamento da
capacidade. Se ρ \textless{} 1, a atividade entra na região de alerta
antes do limite rígido.

Isso cria uma região residual:

\begin{quote}
\textbf{Q \textless{} execução ≤ H}
\end{quote}

Atingir Q não interrompe imediatamente a execução. Q é um \textbf{limiar
informacional}, enquanto H continua sendo o limite de execução.

A semântica de Q possui dois momentos distintos:

\begin{enumerate}
\def\labelenumi{\arabic{enumi}.}
\tightlist
\item
  \textbf{alerta intraday:} ao alcançar Q, a infraestrutura informa o
  emissor e registra a ultrapassagem de forma auditável;
\item
  \textbf{avaliação no corte:} no encerramento da sessão ou em outro
  corte público definido pela especificação, os valores consolidados são
  comparados com Q para determinar a obrigação de disclosure.
\end{enumerate}

Além dessa fronteira quantitativa existe um \textbf{trigger temporal T},
detalhado na Parte III. T é a duração máxima de W, a janela regulatória
iniciada pela primeira execução. Q e T cumprem funções complementares e
podem encerrar W: Q reage à intensidade da utilização; T impede que
atividade abaixo do limiar permaneça indefinidamente sem disclosure.

Essa separação é necessária porque Gross e Net possuem propriedades
temporais diferentes. Gross é cumulativo dentro da janela: BUY e SELL
aumentam Gross, e uma operação oposta não desfaz a ultrapassagem de
\(Q_{\mathrm{gross}}\). Net representa posição e pode recuar. Se Net+
ultrapassa \(Q_{\mathrm{net}+}\) e uma venda posterior reduz a exposição
para abaixo do limiar antes do corte, não há disclosure quantitativo por
Net+ naquele fechamento. O alerta intraday permanece no registro de
auditoria, mas não é convertido retroativamente em obrigação pública.

\begin{quote}
\textbf{Ultrapassar Q net durante a sessão gera alerta; terminar acima
de Q net no corte gera disclosure.}
\end{quote}

\begin{quote}
\textbf{Ultrapassar Q gross não é reversível por operação oposta, porque
Gross acumula as duas pernas.}
\end{quote}

A distinção possui uma função anti-avoidance importante. Se Q fosse
sempre igual a H, a companhia poderia encontrar valor em permanecer
repetidamente em H − ε e esperar apenas o trigger temporal. Com Q abaixo
de H, evitar a avaliação quantitativa exige abrir mão de parcela
material da capacidade disponível. A reversibilidade de Net não apaga a
intervenção necessária para produzir a ida e a volta, porque ambas
continuam consumindo Gross.

Assim, ρ controla quanto da capacidade pode ser utilizada antes que a
atividade entre na região de alerta e possa produzir obrigação
quantitativa no corte.

Nenhum valor específico de ρ é proposto nesta etapa.

Quando o disclosure é devido, ele não deve fingir que a execução foi
exatamente Q nem limitar-se a dizer que ``Q foi ultrapassado''. Q define
a fronteira de avaliação. A quantidade informacionalmente relevante é a
\textbf{execução consolidada efetiva no momento de corte}, da qual
deriva U. Assim, uma especificação-base deve distinguir: a ultrapassagem
intraday gera alerta; a condição consolidada no corte determina a
obrigação; e o disclosure informa quanto foi efetivamente executado.

\begin{center}\rule{0.5\linewidth}{0.5pt}\end{center}

\hypertarget{11-u-utilizacao-como-intensidade-observavel-do-sinal}{}

\hypertarget{11-u-utilizauxe7uxe3o-como-intensidade-observuxe1vel-do-sinal}{%
\subsection{11. U: utilização como intensidade observável do
sinal}\label{11-u-utilizauxe7uxe3o-como-intensidade-observuxe1vel-do-sinal}}

Q define \textbf{quando alertar e em que fronteira avaliar} o
disclosure. U ajuda a representar \textbf{quanto da capacidade foi
efetivamente utilizada}.

Para cada dimensão \(j\):

\begin{quote}
\textbf{\(U_j = \dfrac{\text{execução relevante}_j}{H_j}\)}
\end{quote}

Logo, U varia de 0 a 1 dentro da capacidade permitida.

No lado posicional, é útil separar direções:

\begin{quote}
\textbf{\(U_{\mathrm{net}+}\) = Net positivo / \(H_{\mathrm{net}+}\)}
\end{quote}

\begin{quote}
\textbf{\(U_{\mathrm{net}-}\) = \textbar Net negativo\textbar{} /
\(H_{\mathrm{net}-}\)}
\end{quote}

U não é ``informação'' em si. É uma medida normalizada da intensidade da
decisão observável do emissor em relação à capacidade que ele tinha
disponível. Em uma representação simplificada, a execução observável
produz U; uma função de atualização interpreta essa intensidade
juntamente com direção, tempo e demais características da execução; e o
resultado altera a memória informacional I. Portanto, U é um input da
atualização e não precisa ser o multiplicador direto de informação.

O estado informacional futuro deve ter origem em comportamento
observável. Como a execução do emissor é identificável no regime
especial, a atualização pode utilizar a trajetória executada consolidada
mesmo quando ela não produz disclosure quantitativo adicional por Q. O
disclosure formal agrega, certifica e devolve informação segundo o
protocolo; não é a única forma pela qual a atuação se torna observável.
A cadeia é:

\begin{quote}
H disponível → execução identificável → U observado → eventual
disclosure formal → atualização de informação → decay ao longo do tempo
→ H futuro
\end{quote}

Uma companhia que utiliza 5\% de sua capacidade e uma companhia que
utiliza 95\% não precisam produzir o mesmo sinal, mesmo que ambas tenham
negociado na mesma direção.

A função que transforma U em atualização informacional é experimental.
Ela pode ser linear, côncava, convexa, saturante ou depender de outras
características da execução. O protocolo apenas exige que a hipótese
seja declarada e testável.

No caso de churn, Gross pode ser elevado enquanto Net é próximo de zero.
Isso significa que \(U_{\mathrm{gross}}\) pode ser alto e
\(U_{\mathrm{net}}\) baixo. Essa diferença é precisamente o que permite
separar informação sobre \textbf{atividade} de informação sobre
\textbf{comprometimento posicional}.

\begin{center}\rule{0.5\linewidth}{0.5pt}\end{center}

\hypertarget{12-informacao-produzida-pelo-emissor-informacao-do-mercado-e-atividade-gross}{}

\hypertarget{12-informauxe7uxe3o-produzida-pelo-emissor-informauxe7uxe3o-do-mercado-e-atividade-gross}{%
\subsection{12. Informação produzida pelo emissor, informação do mercado
e atividade
Gross}\label{12-informauxe7uxe3o-produzida-pelo-emissor-informauxe7uxe3o-do-mercado-e-atividade-gross}}

A palavra ``informação'' pode designar fenômenos economicamente
diferentes. A arquitetura os separa explicitamente.

Em síntese, três objetos podem alimentar o estado futuro.
\textbf{\(I_{\mathrm{issuer}}\)} guarda a memória associada às decisões
observáveis do emissor; \textbf{\(I_{\mathrm{market}}\)} representa
informação econômica já disponível externamente; e
\textbf{\(A_{\mathrm{gross}}\)} pode registrar, de forma opcional, a
atividade bruta atribuível à companhia. U funciona como uma medida de
intensidade utilizada na atualização desses estados, sobretudo de
\(I_{\mathrm{issuer}}\), mas não se confunde com nenhum deles.

\hypertarget{12-1-i-issuer-informacao-produzida-pela-acao-divulgada-do-emissor}{}

\hypertarget{121-i_mathrmissuer-informauxe7uxe3o-revelada-pela-auxe7uxe3o-observuxe1vel-do-emissor}{%
\subsubsection{\texorpdfstring{12.1. \(I_{\mathrm{issuer}}\): informação
revelada pela ação observável do
emissor}{12.1. I\_\{\textbackslash mathrm\{issuer\}\}: informação revelada pela ação observável do emissor}}\label{121-i_mathrmissuer-informauxe7uxe3o-revelada-pela-auxe7uxe3o-observuxe1vel-do-emissor}}

\(I_{\mathrm{issuer}}\) representa o estado informacional associado ao
histórico observável de comprometimento posicional da companhia,
composto pelas execuções identificáveis e pelos disclosures formais
aplicáveis. A expressão ``informação produzida pelo emissor'' deve ser
entendida com cuidado: o protocolo não transforma BUY em afirmação de
subavaliação nem SELL em afirmação de sobreavaliação. O dado produzido é
a própria ação econômica observável; a inferência sobre fundamentos e
valuation é produzida pelo mercado.

Ele pode ser direcional:

\begin{itemize}
\tightlist
\item
  \textbf{\(I_{\mathrm{issuer}+}\)}: memória informacional associada a
  BUY líquido;
\item
  \textbf{\(I_{\mathrm{issuer}-}\)}: memória informacional associada a
  SELL líquido.
\end{itemize}

Sua atualização pode depender de \(U_{\mathrm{net}}\), da mudança
efetiva de Treasury, do tempo e de outras características da execução.

Não se exige que \(I_{\mathrm{issuer}+}\) e \(I_{\mathrm{issuer}-}\) se
cancelem. Uma companhia pode ter histórico informativo relevante nas
duas direções. Alternância continua custosa por Gross e pode ser
limitada por \textbf{saturação}, que contempla a perda de valor marginal
de sinais repetidos na mesma direção, e por \textbf{decay}, que
representa a depreciação do valor informacional passado ao longo do
tempo quando o sinal não é renovado. As formas específicas desses
mecanismos são apresentadas na seção seguinte.

\hypertarget{12-2-i-market-informacao-ja-disponivel-externamente}{}

\hypertarget{122-i_mathrmmarket-informauxe7uxe3o-juxe1-disponuxedvel-externamente}{%
\subsubsection{\texorpdfstring{12.2. \(I_{\mathrm{market}}\): informação
já disponível
externamente}{12.2. I\_\{\textbackslash mathrm\{market\}\}: informação já disponível externamente}}\label{122-i_mathrmmarket-informauxe7uxe3o-juxe1-disponuxedvel-externamente}}

\(I_{\mathrm{market}}\) representa informação sobre o estado econômico
da companhia que o mercado consegue observar \textbf{independentemente
das operações do emissor}.

Exemplos possíveis incluem preços de commodities, câmbio, juros ou
outras externalidades diretamente relacionadas ao negócio.

\(I_{\mathrm{market}}\) não é ``crédito ganho pela companhia''. É uma
medida candidata da quantidade ou qualidade de informação econômica que
já está fora da companhia.

A hipótese é que maior observabilidade externa possa reduzir a
assimetria marginal exclusiva do emissor e, portanto, justificar --- em
alguma especificação --- maior flexibilidade operacional. Essa última
passagem precisa ser testada; não é automática.

\hypertarget{12-3-a-gross-informacao-opcional-derivada-da-atividade-bruta}{}

\hypertarget{123-a_mathrmgross-informauxe7uxe3o-opcional-derivada-da-atividade-bruta}{%
\subsubsection{\texorpdfstring{12.3. \(A_{\mathrm{gross}}\): informação
opcional derivada da atividade
bruta}{12.3. A\_\{\textbackslash mathrm\{gross\}\}: informação opcional derivada da atividade bruta}}\label{123-a_mathrmgross-informauxe7uxe3o-opcional-derivada-da-atividade-bruta}}

Além da memória direcional associada à posição líquida, o laboratório
pode testar uma variável opcional chamada
\textbf{\(A_{\mathrm{gross}}\)}. Ela representa um estado informacional
derivado da atividade Gross divulgada: não diz que a companhia aumentou
ou reduziu sua exposição, apenas registra que houve atividade bruta
atribuível ao emissor.

A lógica é:

\begin{quote}
atividade Gross divulgada → possível informação sobre comportamento →
\(A_{\mathrm{gross}}\)
\end{quote}

Mas o regulador pode decidir que puro churn não merece capacidade futura
alguma. Nesse caso:

\begin{quote}
\textbf{\(A_{\mathrm{gross}}\) = 0}
\end{quote}

Outra especificação pode permitir que \(A_{\mathrm{gross}}\) afete
apenas \(H_{\mathrm{gross}}\), e não \(H_{\mathrm{net}}\). Isso é
particularmente atraente porque mantém separado o valor informacional de
``a companhia esteve muito ativa'' do valor informacional de ``a
companhia assumiu posição líquida''.

A arquitetura não canoniza \(A_{\mathrm{gross}}\). Ela apenas deixa
explícita a hipótese para que possa ser incluída ou removida.

\begin{center}\rule{0.5\linewidth}{0.5pt}\end{center}

\hypertarget{13-decay-memoria-e-saturacao}{}

\hypertarget{13-decay-memuxf3ria-e-saturauxe7uxe3o}{%
\subsection{13. Decay, memória e
saturação}\label{13-decay-memuxf3ria-e-saturauxe7uxe3o}}

Decay ocorre sobre \textbf{informação}, não sobre H.

Se nenhum novo sinal relevante surgir, a informação produzida por uma
operação antiga tende a perder capacidade explicativa sobre o estado
econômico atual da companhia.

Assim:

\begin{quote}
\(I_{\mathrm{issuer}}\) elevado hoje → tempo sem novo sinal →
\(I_{\mathrm{issuer}}\) converge gradualmente ao estado neutro
\end{quote}

H muda posteriormente porque é recalculado a partir de um estado
informacional diferente. Não existe necessidade de escrever ``H decai''.

A forma temporal do decay é desconhecida. Pode ser exponencial,
hiperbólica, logística, piecewise ou empiricamente estimada. A tese
exige apenas uma propriedade mínima: sem nova evidência, a memória
informacional não deve permanecer máxima indefinidamente.

Saturação responde a outra pergunta: \textbf{qual é o limite do efeito
que informação acumulada pode exercer sobre H?}

Se H₀ é a referência neutra e \(H_{\max}\) o teto admissível, a
informação deve mapear H para algum ponto dentro desse intervalo. A
função pode possuir retornos marginais decrescentes, mas isso é uma
propriedade candidata, não uma definição necessária.

Essa separação produz quatro objetos distintos:

\begin{itemize}
\tightlist
\item
  H₀: referência estrutural;
\item
  \(H_{\max}\): teto de capacidade sob a especificação;
\item
  I: estado informacional;
\item
  decay: envelhecimento de I.
\end{itemize}

Ela evita que ``memória'', ``capacidade'' e ``saturação'' sejam tratados
como sinônimos.

\begin{center}\rule{0.5\linewidth}{0.5pt}\end{center}

\hypertarget{iii-tempo-disclosure-e-transicao-de-estado}{}

\hypertarget{iii-tempo-disclosure-e-transiuxe7uxe3o-de-estado-1}{%
\section{III. Tempo, disclosure e transição de
estado}\label{iii-tempo-disclosure-e-transiuxe7uxe3o-de-estado-1}}

\hypertarget{14-w-janela-regulatoria-disclosure-e-transicao-de-estado}{}

\hypertarget{14-w-janela-regulatuxf3ria-disclosure-e-transiuxe7uxe3o-de-estado}{%
\subsection{14. W: janela regulatória, disclosure e transição de
estado}\label{14-w-janela-regulatuxf3ria-disclosure-e-transiuxe7uxe3o-de-estado}}

O tempo do protocolo é organizado por \textbf{W}, a janela regulatória
dentro da qual a atuação do emissor é acumulada e avaliada. W não é
apenas um intervalo de calendário: é o universo comum ao qual pertencem
H, Q e T.

\begin{itemize}
\tightlist
\item
  \textbf{H pertence a W:} a capacidade é calculada na abertura da
  janela, permanece travada durante sua vigência e é consumida pelas
  execuções nela realizadas;
\item
  \textbf{Q pertence a W:} o limiar quantitativo é definido sobre os H
  daquela janela e determina quando a utilização entra na região de
  alerta e de avaliação de disclosure;
\item
  \textbf{T pertence a W:} o prazo começa com a abertura da janela,
  limita sua duração e funciona como trigger temporal de disclosure caso
  Q não a encerre antes.
\end{itemize}

No baseline, a primeira execução elegível do emissor abre \(W_n\). Antes
dela não existe capacidade temporal sendo consumida nem relógio T em
curso. A partir dela, todas as execuções elegíveis do mesmo emissor são
atribuídas a \(W_n\) até seu encerramento.

O protocolo possui, portanto, dois gatilhos complementares de devolução
de informação. \textbf{Q} reage à intensidade da utilização; \textbf{T}
impede que uma janela permaneça aberta indefinidamente com atividade
abaixo de Q. Se o corte aplicável confirmar a condição de Q, o
disclosure quantitativo é devido e fecha W. Se Q não produzir esse
encerramento, o vencimento de T aciona o disclosure temporal e também
fecha W.

A ultrapassagem intraday de Q continua sendo alerta, não encerramento
automático. Para \(Q_{\mathrm{gross}}\), a ultrapassagem persiste até o
corte porque Gross não diminui. Para \(Q_{\mathrm{net}+}\) e
\(Q_{\mathrm{net}-}\), uma operação oposta pode devolver a exposição
para baixo do limiar antes da consolidação. Nesse caso, não há
fechamento de W por Q, embora a ultrapassagem e sua reversão permaneçam
no registro auditável.

Quando Q é confirmado, a execução não precisa ter parado exatamente no
limiar. A companhia pode ter utilizado a capacidade residual até H antes
do corte, mas Q não cria capacidade nova nem recicla a parcela
consumida. O disclosure informa \textbf{o que foi efetivamente executado
em W} --- direção, quantidades, utilização e demais campos do protocolo
---, e não uma quantidade ficticiamente igual a Q.

O fechamento da sessão é o corte público do baseline. Ele separa o
alerta operacional da obrigação pública, incorpora a trajetória completa
de Net e impede que o disclosure reabra capacidade no mesmo fluxo
intraday. Encerrada \(W_n\), o estado é atualizado; a capacidade não
utilizada expira; e \(W_{n+1}\) pode começar \textbf{no pregão
seguinte}, nunca no mesmo pregão do disclosure que fechou a janela
anterior.

Em forma compacta:

\begin{quote}
primeira execução → abertura de \(W_n\) e início de T → consumo de H e
monitoramento de Q → corte → fechamento por Q ou T → disclosure →
atualização do estado → elegibilidade de \(W_{n+1}\) no pregão seguinte
\end{quote}

Não existem janelas sobrepostas para o mesmo emissor.

\begin{center}\rule{0.5\linewidth}{0.5pt}\end{center}

\hypertarget{15-modelo-temporal-rolling}{}

\hypertarget{15-modelo-temporal-rolling}{%
\subsection{15. Modelo temporal
rolling}\label{15-modelo-temporal-rolling}}

O modelo \textbf{rolling} é a expressão mais direta desse ciclo. A
primeira execução elegível abre \(W_n\) e inicia a contagem de T naquele
instante. Se Q não fechar a janela antes, T vence após a duração
especificada, independentemente de fronteiras comuns do calendário.

Esse desenho faz o relógio acompanhar a atividade efetiva do emissor.
Uma companhia inativa não mantém W artificialmente aberta; companhias
que começam a operar em momentos diferentes possuem janelas distintas. O
mercado observa as execuções e os disclosures do protocolo sem precisar
tratar a abertura administrativa de W como um evento econômico separado.

Após o encerramento:

\begin{quote}
\(W_n\) → disclosure por Q ou T → atualização de estado → cálculo de
\(H_{n+1}\) → \(W_{n+1}\) elegível no pregão seguinte
\end{quote}

Q continua necessário mesmo com T. Sem a fronteira quantitativa, o
emissor poderia consumir parcela muito elevada de H logo após a abertura
e devolver informação apenas no prazo máximo. Q antecipa o encerramento
quando a intensidade da utilização assim exige.

\begin{center}\rule{0.5\linewidth}{0.5pt}\end{center}

\hypertarget{16-modelo-temporal-calendarizado}{}

\hypertarget{16-modelo-temporal-calendarizado}{%
\subsection{16. Modelo temporal
calendarizado}\label{16-modelo-temporal-calendarizado}}

O modelo \textbf{calendarizado} mantém a primeira execução como ato que
ativa W, mas ancora seu limite externo em um ciclo público previamente
conhecido --- por exemplo, semana, quinzena ou mês regulatório. Se
nenhuma operação ocorrer, não existe W ativa; se a companhia operar, a
janela se encerra no primeiro evento entre a confirmação de Q, o
vencimento de T e a fronteira calendarizada aplicável.

A vantagem é a previsibilidade: mercado e emissor conhecem
antecipadamente o limite máximo do ciclo. O custo potencial é que essa
fronteira se torna informação estratégica. Participantes podem adaptar
liquidez, ordens e antecipação à proximidade do corte. Sazonalidade
microestrutural em torno dessas datas é, portanto, uma hipótese a ser
testada.

Como no rolling, H permanece fixo enquanto W está aberta, a execução
consome capacidade e nenhum disclosure recicla H dentro da mesma janela.

\begin{center}\rule{0.5\linewidth}{0.5pt}\end{center}

\hypertarget{17-o-que-realmente-diferencia-as-duas-branches}{}

\hypertarget{17-duas-topologias-temporais-para-o-mesmo-mecanismo}{%
\subsection{17. Duas topologias temporais para o mesmo
mecanismo}\label{17-duas-topologias-temporais-para-o-mesmo-mecanismo}}

Os modelos calendarizado e rolling implementam o mesmo motor econômico:

\begin{quote}
estado → primeira execução e abertura de W → cálculo e trava de H →
execução → U → alerta Q → fechamento por Q ou T → disclosure →
atualização de informação → novo estado → novo H
\end{quote}

A diferença é \textbf{como o tempo delimita uma janela regulatória}.

Nos dois modelos, a primeira execução ativa W e Q pode encerrá-la. A
diferença está na âncora temporal remanescente. No rolling, T é contado
exclusivamente a partir da primeira execução. No calendarizado, W também
se submete a uma fronteira pública externa, que pode antecipar seu
encerramento.

Essa separação é útil porque permite comparar as duas topologias
mantendo constantes função de H, parâmetros, agentes e condições de
mercado. O modelo calendarizado oferece fronteiras temporais simples e
públicas; o rolling reduz a dependência de um relógio comum, mas
introduz transições endógenas associadas à própria atividade do emissor.
Nenhum dos dois é tratado como superior a priori.

\begin{center}\rule{0.5\linewidth}{0.5pt}\end{center}

\hypertarget{18-cooldown-reversao-e-fato-relevante}{}

\hypertarget{18-intervalos-de-reauxe7uxe3o-reversuxe3o-e-fato-relevante}{%
\subsection{18. Intervalos de reação, reversão e Fato
Relevante}\label{18-intervalos-de-reauxe7uxe3o-reversuxe3o-e-fato-relevante}}

Além das janelas W e do prazo T, o desenho pode testar a necessidade de
um \textbf{intervalo adicional sem nova intervenção do emissor} após
determinado disclosure. Chamemos esse intervalo opcional de \textbf{C}.

C existe para responder a uma pergunta empírica: a fronteira temporal já
criada pelo disclosure e pela próxima sessão oferece tempo suficiente
para que o mercado reaja, ou é necessário reservar um período adicional?
Se a própria separação entre sessões for suficiente, C pode ser zero. Se
não for, valores positivos podem ser testados. A tese não presume
antecipadamente que esse intervalo adicional seja necessário.

A reversão de direção é um problema relacionado, mas distinto. Uma
companhia pode comprar e posteriormente vender ações em Treasury por
mudança de avaliação, necessidade de caixa, realização de valor ou outra
razão econômica legítima. Ao mesmo tempo, reversões muito rápidas podem
ser utilizadas para tentar monetizar a reação produzida pela própria
atuação identificada.

A primeira proteção é estrutural: \textbf{reversão não restitui Gross}.
Comprar e depois vender pode reduzir Net, mas consome intervenção nas
duas pernas. O laboratório pode ainda testar uma fricção direcional
específica, como degradação parcial da capacidade futura após reversão
ou um intervalo mínimo antes da direção oposta. Essas alternativas
pertencem ao espaço experimental porque uma trava excessiva também cria
risco: terceiros podem explorar o fato de que o emissor ficou
previsivelmente preso a uma direção.

Fato Relevante permanece externo à máquina do protocolo. Sua divulgação
pode alterar radicalmente a interpretação das operações anteriores e o
valor da informação privada remanescente, mas não produz automaticamente
bônus, reset ou novo H. As obrigações jurídicas de Fato Relevante
continuam existindo por sua própria lógica
(\protect\hyperlink{ref-2}{CVM, 2021}).

\begin{center}\rule{0.5\linewidth}{0.5pt}\end{center}

\hypertarget{iv-informacao-aprendizagem-e-integridade}{}

\hypertarget{iv-informauxe7uxe3o-aprendizagem-e-integridade-1}{%
\section{IV. Informação, aprendizagem e
integridade}\label{iv-informauxe7uxe3o-aprendizagem-e-integridade-1}}

\hypertarget{19-a-companhia-como-trader-identificavel-e-potencialmente-informado}{}

\hypertarget{19-a-companhia-como-trader-identificuxe1vel-e-potencialmente-informado}{%
\subsection{19. A companhia como trader identificável e potencialmente
informado}\label{19-a-companhia-como-trader-identificuxe1vel-e-potencialmente-informado}}

A proposta depende de uma ideia já central à teoria de microestrutura:
\textbf{ordens e transações podem carregar informação}.

Glosten e Milgrom mostram, em um modelo clássico de market making com
traders heterogeneamente informados, que a presença de traders com
informação superior gera componente de seleção adversa no spread e que
preços de transação transmitem informação. Hasbrouck trata empiricamente
o efeito informacional de trades por meio de seu impacto permanente
sobre preços e encontra relações entre tamanho do trade, impacto e
spread (\protect\hyperlink{ref-4}{Glosten e Milgrom, 1985};
\protect\hyperlink{ref-5}{Hasbrouck, 1991}).

Esta tese adiciona uma característica institucional incomum: o trader
potencialmente informado não é apenas inferido estatisticamente.
\textbf{Sua identidade como emissor é pública.}

Isso pode tornar a ação da companhia um sinal particularmente forte.

O protocolo não revela a informação privada subjacente. Ele revela
comportamento econômico padronizado: direção, intensidade, utilização da
capacidade, timing e histórico.

Investidores sofisticados já dedicam recursos à inferência de informação
a partir de fluxo, trades e mudanças de quotes. O mecanismo proposto não
cria do zero essa atividade. Ele altera a qualidade e a atribuição de
uma fonte específica de fluxo ao tornar pública a identidade do emissor.

\begin{quote}
\textbf{O dado pode ser democratizado; a capacidade de interpretá-lo
continua competitiva.}
\end{quote}

\begin{center}\rule{0.5\linewidth}{0.5pt}\end{center}

\hypertarget{20-reputacao-antecipacao-e-reflexividade}{}

\hypertarget{20-reputauxe7uxe3o-antecipauxe7uxe3o-e-reflexividade}{%
\subsection{20. Reputação, antecipação e
reflexividade}\label{20-reputauxe7uxe3o-antecipauxe7uxe3o-e-reflexividade}}

O protocolo não converte reputação em pontuação regulatória. Reputação é
uma crença endógena do mercado formada a partir do histórico público:

\begin{quote}
emissor X comprou → quanto comprou → quanto de H utilizou → o que
aconteceu depois → como se comportou em reversões → como o mercado
reagiu
\end{quote}

Se BUY de determinada companhia historicamente precede informação
econômica positiva, participantes podem reagir mais rapidamente a novos
BUYs. Se a companhia alterna repetidamente entre BUY e SELL após reações
fortes, o mercado também pode aprender a descontar, contrariar ou punir
esse padrão. Isso não é, por si só, abuso. É aprendizagem.

A companhia também observa o mercado. O protocolo contém, portanto, dois
pares simultâneos de ação e reação:

\begin{quote}
\textbf{companhia → mercado:} execução, sinal observável, reação de
preço e liquidez
\end{quote}

\begin{quote}
\textbf{mercado → companhia:} antecipação, arbitragem, spreads, reação
ao histórico e alteração do custo de entrada ou saída
\end{quote}

O desenho não deve presumir cooperação de nenhum dos lados. Uma
companhia racional procura a melhor política permitida pelo protocolo;
participantes racionais procuram explorar tanto a informação da
companhia quanto as restrições que o próprio protocolo lhe impõe. Uma
regra que torne reversão excessivamente difícil, por exemplo, pode
permitir que terceiros negociem sabendo que o emissor está
regulatoriamente encurralado.

Por isso, comportamento estratégico não é uma anomalia a ser removida do
modelo. É parte do objeto. O regulador deve perguntar como a companhia
poderia explorar mercado e protocolo e, simetricamente, como o mercado
poderia sufocar ou explorar a companhia.

\begin{center}\rule{0.5\linewidth}{0.5pt}\end{center}

\hypertarget{21-o-risco-de-o-sinal-quebrar-o-mercado}{}

\hypertarget{21-saliuxeancia-familiaridade-e-aprendizagem-do-sinal}{%
\subsection{21. Saliência, familiaridade e aprendizagem do
sinal}\label{21-saliuxeancia-familiaridade-e-aprendizagem-do-sinal}}

O protocolo cria deliberadamente um sinal potencialmente poderoso:

\begin{quote}
\textbf{trader identificado + provável vantagem informacional + direção
+ intensidade de utilização da capacidade}
\end{quote}

A intensidade desse sinal pode depender não apenas de U e da reputação
do emissor, mas também da \textbf{raridade do próprio evento}. Uma
primeira operação identificada sob um regime novo pode produzir
estardalhaço, overshooting e proteção agressiva de market makers. O
mesmo U, depois de dezenas de interações observadas, pode receber
interpretação muito diferente. Familiaridade institucional e
familiaridade operacional são, portanto, dimensões distintas.

Isso não transforma choque informacional em defeito por definição.
Saliência não é manipulação; volatilidade não é necessariamente falha;
overreaction inicial não implica desequilíbrio permanente. Uma
trajetória plausível é:

\begin{quote}
operação identificada rara → reação excessiva → companhia e mercado
observam o resultado → agentes recalibram → novas operações recebem
interpretação menos mecânica → possível estabilização
\end{quote}

Também é possível que a estabilização não ocorra. Market makers podem
reduzir liquidez, seguidores podem amplificar excessivamente o sinal,
reputação pode gerar múltiplos equilíbrios ou a identificação do emissor
pode dominar a produção independente de informação. A existência de
aprendizagem saudável é hipótese de validação, não premissa.

Os mecanismos apresentados na seção 19 tornam ambos os resultados
plausíveis: a identificação de um trader potencialmente informado pode
ampliar proteção contra seleção adversa, enquanto sua presença também
pode fornecer contraparte e liquidez. A evidência específica sobre
recompras, discutida na seção 29, encontra sinais diferentes conforme
mercado e amostra. O efeito líquido não pode ser presumido pela
arquitetura.

O protocolo não pretende escolher a reação correta do mercado.
Arbitradores, market makers, especuladores e investidores podem copiar,
contrariar ou ignorar o emissor. O objetivo é preservar condições para
que essas respostas existam e sejam aprendidas.

Grossman e Stiglitz mostram, em contexto mais geral, por que mercados
perfeitamente informacionalmente eficientes são incompatíveis com
incentivos para aquisição custosa de informação
(\protect\hyperlink{ref-6}{Grossman e Stiglitz, 1980}). Esta tese não
aplica o modelo deles mecanicamente ao protocolo; utiliza-o como
lembrete de que \textbf{mais informação pública não é automaticamente
equivalente a uma estrutura informacional melhor}.

A questão de viabilidade é dinâmica: a saliência inicial converge para
um processo informacional administrável ou destrói liquidez,
executabilidade ou descoberta descentralizada de preço antes que o
aprendizado ocorra?

\begin{center}\rule{0.5\linewidth}{0.5pt}\end{center}

\hypertarget{22-manipulacao-gross-e-transferencia-de-valor-a-terceiros}{}

\hypertarget{22-captura-parcial-de-valor-fabricauxe7uxe3o-de-sinal-e-fronteira-da-manipulauxe7uxe3o}{%
\subsection{22. Captura parcial de valor, fabricação de sinal e
fronteira da
manipulação}\label{22-captura-parcial-de-valor-fabricauxe7uxe3o-de-sinal-e-fronteira-da-manipulauxe7uxe3o}}

O dado divulgado pelo protocolo é ``a companhia comprou'' ou ``a
companhia vendeu'', acompanhado de quantidade e utilização de
capacidade. Ele \textbf{não} é ``a companhia declarou que está barato''
ou ``a companhia declarou que está caro''. O valuation permanece uma
inferência dos participantes.

A companhia continua podendo capturar parte da alta que ajudou a
antecipar. Ela compra enquanto o preço ainda não incorporou
integralmente sua informação; se estiver correta, a posição em Treasury
se valoriza. O protocolo não elimina esse retorno --- eliminá-lo
retiraria parte substancial do incentivo econômico para participar do
regime.

O que o protocolo procura impedir é a captura \textbf{integral e
unilateral} do movimento. Para aumentar a posição, a companhia consome
H; ao utilizar parcela relevante, aproxima-se de Q e antecipa o
disclosure; ao atingir Q ou T, fecha W e precisa aguardar a próxima
janela; e, uma vez identificada, sua atuação pode alterar preço e
liquidez contra as execuções seguintes. Quanto maior a tentativa de
monetizar a vantagem, maior a parcela da própria informação devolvida ao
mercado e menor o espaço para continuar executando antes da reação. Essa
é uma propriedade autorregulatória pretendida da arquitetura, não uma
garantia já demonstrada.

Essa estrutura permite separar duas situações. Na primeira, a companhia
assume posição econômica genuína, o mercado participa do repricing e, em
estado posterior de preço e informação, a companhia realiza parte de
Treasury. O ganho corporativo pode coexistir com participação de
terceiros na descoberta de preço. Se a decisão estiver errada, o
patrimônio da própria companhia absorve a perda e o histórico público
pode reduzir a credibilidade de sinais futuros.

Na segunda situação, a finalidade econômica predominante da operação é
fabricar percepção artificial de demanda, convicção ou valor para
induzir terceiros a negociar e monetizar a reação produzida. A execução
pode ser real e ainda assim levantar problema de manipulação. H, Q, T e
a identificação do emissor não tornam intenção manipulativa impossível;
procuram limitar sua escala, elevar seu custo e tornar a trajetória
reconstruível.

A arquitetura não pretende resolver intenção por álgebra. Gross impede
que BUY e SELL sucessivos apaguem a intervenção realizada; H limita
escala em W; Q e T encerram a janela e disciplinam a devolução
informacional; Treasury limita SELL físico; reversões permanecem
observáveis; e especificações de H podem tornar persistência econômica
mais capaz de gerar escala do que alternância estratégica.

Isso conduz a uma premissa de elaboração do sistema:

\begin{quote}
\textbf{Nenhuma estratégia exploratória deve conseguir escala relevante
sem consumir capacidade, assumir risco econômico, deixar rastro
observável ou destruir a própria capacidade futura.}
\end{quote}

A premissa é um alvo de convergência, não uma promessa de
impossibilidade de exploit. Se uma estratégia barata, escalável e
repetível permite extrair valor porque o protocolo garante mecanicamente
a reação de terceiros, existe uma fragilidade estrutural. Se a
exploração exige capital, Gross, tempo, risco de mercado e perda de
opcionalidade futura, a própria dinâmica pode discipliná-la.

Essa disciplina pode vir também do mercado. Uma companhia que
repetidamente tente explorar BUY → valorização → SELL pode ensinar
participantes a não seguir o próximo BUY, a operar contra ele ou a
alterar spreads e profundidade. A exploração pode, assim, reduzir a
eficácia futura da própria exploração. A engenharia deve
\textbf{fortalecer}, e não substituir, esse mecanismo endógeno sempre
que possível.

Também importa observar terceiros. Participantes que inferem
publicamente o sinal competem informacionalmente; participantes que
antecipam ordens futuras por acesso indevido apresentam problema
diferente, potencialmente envolvendo vazamento ou front running.
Investigações sérias podem exigir dados não públicos de ordens,
cancelamentos, contrapartes e lead-lag.

\begin{center}\rule{0.5\linewidth}{0.5pt}\end{center}

\hypertarget{23-disciplina-endogena-e-enforcement-ex-post}{}

\hypertarget{23-disciplina-enduxf3gena-desenho-adversarial-e-enforcement-ex-post}{%
\subsection{23. Disciplina endógena, desenho adversarial e enforcement
ex
post}\label{23-disciplina-enduxf3gena-desenho-adversarial-e-enforcement-ex-post}}

Parte das propriedades desejadas pode ser imposta mecanicamente:

\begin{quote}
Gross não pode ultrapassar \(H_{\mathrm{gross}}\).
\end{quote}

\begin{quote}
Net positivo não pode ultrapassar \(H_{\mathrm{net}+}\).
\end{quote}

\begin{quote}
Net negativo não pode ultrapassar \(H_{\mathrm{net}-}\).
\end{quote}

\begin{quote}
SELL não pode ultrapassar Treasury disponível.
\end{quote}

\begin{quote}
Operação oposta não apaga Gross consumido.
\end{quote}

Essas restrições não são punições. São invariantes do regime.

A primeira forma de disciplina é endógena:

\begin{quote}
uso da capacidade → consumo de H → alerta Q → fechamento de W por Q ou T
→ disclosure → reação e aprendizagem do mercado
\end{quote}

A segunda é jurídica e supervisória:

\begin{quote}
conduta → investigação → interpretação → eventual sanção
\end{quote}

A primeira não substitui a segunda. Manipulação, fraude, coordenação,
vazamento, informação falsa e estratégias não antecipadas continuam
exigindo julgamento. Da mesma forma, conformidade mecânica com H, Q e T
não deve funcionar como escudo para uma conduta cuja natureza econômica
seja manipulativa.

Dois princípios de desenho seguem dessa separação:

\begin{quote}
\textbf{Invariante quando possível; supervisão quando necessária;
julgamento quando inevitável.}
\end{quote}

\begin{quote}
\textbf{Sempre que possível, a arquitetura deve fortalecer mecanismos
endógenos de disciplina em vez de substituí-los por proibições ex ante.}
\end{quote}

O protocolo deve ser desenhado contra o comportamento racional dos dois
lados, e não a partir da expectativa de comportamento cooperativo. O
teste adversarial relevante pergunta simultaneamente qual seria a melhor
exploração disponível à companhia e qual seria a melhor resposta do
mercado às ações e restrições da companhia.

\begin{center}\rule{0.5\linewidth}{0.5pt}\end{center}

\hypertarget{v-especificacao-calibracao-e-laboratorio-regulatorio}{}

\hypertarget{v-especificauxe7uxe3o-calibrauxe7uxe3o-e-laboratuxf3rio-regulatuxf3rio-1}{%
\section{V. Especificação, calibração e laboratório
regulatório}\label{v-especificauxe7uxe3o-calibrauxe7uxe3o-e-laboratuxf3rio-regulatuxf3rio-1}}

\hypertarget{24-arquitetura-especificacao-e-calibracao}{}

\hypertarget{24-arquitetura-especificauxe7uxe3o-e-calibrauxe7uxe3o}{%
\subsection{24. Arquitetura, especificação e
calibração}\label{24-arquitetura-especificauxe7uxe3o-e-calibrauxe7uxe3o}}

A tese distingue três camadas.

\hypertarget{arquitetura}{%
\subsubsection{Arquitetura}\label{arquitetura}}

Define o significado dos objetos:

\begin{itemize}
\tightlist
\item
  W como janela regulatória à qual pertencem H, Q e T;
\item
  Gross e Net;
\item
  \(H_{\mathrm{gross}}\), \(H_{\mathrm{net}+}\) e \(H_{\mathrm{net}-}\);
\item
  Q como limiar de alerta e avaliação no corte;
\item
  U como utilização relativa;
\item
  Treasury;
\item
  abertura de W pela primeira execução, encerramento por Q ou T e
  transição para a janela seguinte;
\item
  informação e decay;
\item
  topologia Calendar ou Rolling.
\end{itemize}

\hypertarget{especificauxe7uxe3o}{%
\subsubsection{Especificação}\label{especificauxe7uxe3o}}

Escolhe como os objetos se relacionam:

\begin{itemize}
\tightlist
\item
  quais variáveis entram em H;
\item
  como U atualiza \(I_{\mathrm{issuer}}\);
\item
  se existe \(A_{\mathrm{gross}}\);
\item
  como \(I_{\mathrm{market}}\) é representado;
\item
  qual função de saturação é usada;
\item
  como decay funciona;
\item
  como H₀ e \(H_{\max}\) são definidos;
\item
  se plano público voluntário produz bônus moderado de capacidade
  inicial;
\item
  como persistência e reversão afetam a progressão direcional de H.
\end{itemize}

\hypertarget{calibrauxe7uxe3o}{%
\subsubsection{Calibração}\label{calibrauxe7uxe3o}}

Escolhe ou estima valores:

\begin{itemize}
\tightlist
\item
  tamanho de H₀;
\item
  \(H_{\max}\);
\item
  ρ = Q/H;
\item
  duração T;
\item
  parâmetros de decay;
\item
  sensibilidade de H à informação;
\item
  neutralidade de Treasury;
\item
  pesos ou funções das variáveis externas;
\item
  bônus por plano público, se existir;
\item
  aceleração de H por persistência;
\item
  degradação de capacidade por reversão, se existir.
\end{itemize}

A arquitetura pode estar madura enquanto especificação e calibração
permanecem abertas. Isso não é deficiência; é condição para que a
hipótese possa ser testada sem inventar precisão.

O laboratório acrescenta uma quarta distinção metodológica, sem alterar
essas três camadas. Um \textbf{cenário} descreve uma trajetória exógena
de ordens associada a uma companhia e a uma estratégia. Uma
\textbf{especificação} contém as relações formais de memória, capacidade
e intrawindow. Uma \textbf{avaliação} aplica uma especificação a um
cenário a partir de um snapshot de dados:

\begin{quote}
\textbf{resultado = F(snapshot de dados, cenário, especificação)}
\end{quote}

Cada avaliação elementar corresponde a uma combinação especificação ×
cenário. Comparações agregam essas avaliações de duas formas
controladas: várias especificações sobre um cenário fixo, ou vários
cenários sob uma especificação fixa. Essa separação permite alterar a
estratégia de ordens mantendo a regra constante, ou alterar a regra
mantendo a trajetória solicitada constante. O cenário não contém Q, H ou
disclosure; esses objetos são resolvidos pela especificação e pelo
estado da companhia.

\begin{center}\rule{0.5\linewidth}{0.5pt}\end{center}

\hypertarget{25-observabilidade-externa-e-flexibilidade-operacional}{}

\hypertarget{25-observabilidade-externa-e-flexibilidade-operacional}{%
\subsection{25. Observabilidade externa e flexibilidade
operacional}\label{25-observabilidade-externa-e-flexibilidade-operacional}}

Esta seção retoma \(I_{\mathrm{market}}\), apresentado na seção 12.2, e
desloca a discussão da ontologia para a especificação. A vantagem
informacional do emissor não precisa ter a mesma intensidade em todos os
ativos ou estados do mercado.

Parte da informação economicamente relevante de uma companhia pode ser
inferida por terceiros a partir de variáveis públicas. Uma empresa
fortemente dependente de commodity cotada, câmbio ou juros pode estar
exposta a choques que o mercado observa em tempo real.

Isso não elimina informação privada. A companhia continua conhecendo
contratos, custos, produtividade, estratégia, hedge, execução e inúmeras
outras variáveis internas.

A hipótese é marginal:

\begin{quote}
\textbf{maior observabilidade externa → potencialmente menor parcela da
informação economicamente relevante exclusiva ao emissor → possibilidade
de flexibilização de H}
\end{quote}

A última seta não deve ser presumida.

Setor pode servir como proxy grosseira quando características melhores
não estiverem disponíveis, mas a preferência metodológica é decompor a
ideia em variáveis objetivas quando possível: exposição a commodity,
sensibilidade cambial, dependência de preços públicos, concentração
geográfica ou outras externalidades mensuráveis.

O setor não recebe capacidade. O ticker recebe H.

\begin{center}\rule{0.5\linewidth}{0.5pt}\end{center}

\hypertarget{26-variaveis-modulares-e-variable-builder}{}

\hypertarget{26-variuxe1veis-modulares-e-variable-builder}{%
\subsection{26. Variáveis modulares e Variable
Builder}\label{26-variuxe1veis-modulares-e-variable-builder}}

A existência de interações plausíveis cuja relevância ainda é
desconhecida cria um problema metodológico: ignorá-las pode empobrecer o
modelo; incorporá-las definitivamente pode transformar especulação em
regra.

A solução proposta é modularidade experimental.

\begin{quote}
\textbf{Quando a análise identifica uma interação economicamente
plausível, mas não fornece evidência suficiente para transformá-la em
regra, o mecanismo pode torná-la uma hipótese modular e testável.}
\end{quote}

O laboratório deve separar:

\begin{quote}
\textbf{dados → variáveis → função de H}
\end{quote}

Dados representam fatos observados e não pertencem à especificação.
Treasury, preço de referência, ações emitidas, ações em circulação e
free float devem ser mapeados semanticamente no snapshot da companhia. A
especificação pode selecionar variáveis derivadas desses dados e definir
como elas entram em H, mas não pode transformar uma posição física
observada em parâmetro discricionário da própria regra.

Uma variável pode vir de dados nativos do mecanismo, dados de mercado ou
fonte externa. Ela pode possuir transformação, normalização, lag,
bounds, frequência e política de missing data explicitamente
documentados.

Exemplos de variáveis primitivas:

\begin{itemize}
\tightlist
\item
  ADTV;
\item
  depth;
\item
  spread;
\item
  free float;
\item
  volatilidade;
\item
  Treasury;
\item
  preço de commodity;
\item
  câmbio;
\item
  existência e dimensão de plano público voluntário.
\end{itemize}

Exemplos derivados:

\begin{itemize}
\tightlist
\item
  utilização U;
\item
  exposição externa estimada;
\item
  estado informacional \(I_{\mathrm{issuer}}\);
\item
  estado de atividade \(A_{\mathrm{gross}}\);
\item
  persistência direcional;
\item
  frequência de reversão;
\item
  familiaridade histórica do mercado com operações identificadas.
\end{itemize}

Cada função de H deve ser uma pipeline independente. O regulador pode
testar, por exemplo:

\begin{quote}
\(H_{\mathrm{net}+}\) = função(H₀, \(I_{\mathrm{issuer}+}\), posição)
\end{quote}

contra:

\begin{quote}
\(H_{\mathrm{net}+}\) = função(H₀, \(I_{\mathrm{issuer}+}\), posição,
\(I_{\mathrm{market}}\))
\end{quote}

ou testar \(H_{\mathrm{gross}}\) com e sem \(A_{\mathrm{gross}}\).

A modularidade não é feature creep. Ela existe para tornar
\textbf{barato rejeitar hipóteses}.

Se uma variável não melhora a estabilidade, a explicabilidade ou o
desempenho do mecanismo, ela deve poder ser removida sem alterar a
ontologia central.

\begin{center}\rule{0.5\linewidth}{0.5pt}\end{center}

\hypertarget{27-invariantes-e-espaco-experimental}{}

\hypertarget{27-invariantes-e-espauxe7o-experimental}{%
\subsection{27. Invariantes e espaço
experimental}\label{27-invariantes-e-espauxe7o-experimental}}

Algumas propriedades pertencem à arquitetura e não devem ser tratadas
como sliders arbitrários:

\begin{longtable}[]{@{}
  >{\raggedright\arraybackslash}p{(\columnwidth - 2\tabcolsep) * \real{0.7600}}
  >{\raggedright\arraybackslash}p{(\columnwidth - 2\tabcolsep) * \real{0.2400}}@{}}
\toprule\noalign{}
\begin{minipage}[b]{\linewidth}\raggedright
Propriedade
\end{minipage} & \begin{minipage}[b]{\linewidth}\raggedright
Status
\end{minipage} \\
\midrule\noalign{}
\endhead
\bottomrule\noalign{}
\endlastfoot
Gross e Net são distintos & estrutural \\
\(H_{\mathrm{gross}}\), \(H_{\mathrm{net}+}\) e \(H_{\mathrm{net}-}\)
são capacidades separadas & estrutural \\
H é hard cap & estrutural \\
ordens acima dos limites têm execução parcial e residual rejeitado &
estrutural \\
Q gera alerta intraday e é avaliado no corte para disclosure &
estrutural \\
Q ≤ H & estrutural \\
SELL ≤ Treasury & estrutural \\
Treasury não pode ficar negativa & estrutural \\
reversão não restitui Gross & estrutural \\
H, Q e T são resolvidos dentro de W & estrutural \\
disclosure confirmado por Q ou acionado por T encerra W & estrutural \\
uma nova W não começa antes do pregão seguinte & baseline estrutural \\
decay atua sobre informação, não diretamente sobre H & estrutural \\
Fato Relevante não reseta automaticamente H & baseline estrutural \\
derivativos sobre a própria ação ficam fora do regime especial &
baseline estrutural \\
\end{longtable}

Outras propriedades pertencem ao espaço experimental:

\begin{itemize}
\tightlist
\item
  valor de Q/H;
\item
  H₀ e \(H_{\max}\);
\item
  função de H;
\item
  função U → \(I_{\mathrm{issuer}}\);
\item
  forma do decay;
\item
  forma da saturação;
\item
  existência de \(A_{\mathrm{gross}}\);
\item
  uso de \(I_{\mathrm{market}}\);
\item
  variáveis externas;
\item
  duração T;
\item
  modelo calendarizado versus modelo rolling;
\item
  eventual intervalo adicional de reação C;
\item
  bônus moderado por plano público;
\item
  velocidade e forma de crescimento de H com execuções sucessivas;
\item
  degradação parcial de H após reversão;
\item
  assimetria de capacidade na direção oposta;
\item
  familiaridade e taxa-base de operações identificadas.
\end{itemize}

Essa separação impede que o laboratório vire um ambiente em que
``qualquer coisa vale''.

\begin{center}\rule{0.5\linewidth}{0.5pt}\end{center}

\hypertarget{28-criterio-de-viabilidade}{}

\hypertarget{28-crituxe9rio-de-viabilidade}{%
\subsection{28. Critério de
viabilidade}\label{28-crituxe9rio-de-viabilidade}}

O protocolo não precisa maximizar uma função escalar única de bem-estar
para ser testável.

É suficiente definir um vetor de resultados e restrições mínimas.

Do lado da companhia, interessam pelo menos:

\begin{itemize}
\tightlist
\item
  capacidade economicamente útil de execução;
\item
  custo de implementação;
\item
  retorno da Treasury;
\item
  flexibilidade de capital allocation;
\item
  previsibilidade e segurança jurídica para operações de boa-fé.
\end{itemize}

Do lado do mercado:

\begin{itemize}
\tightlist
\item
  spread;
\item
  depth;
\item
  volume;
\item
  impacto de preço;
\item
  velocidade e qualidade da descoberta de preço;
\item
  concentração de ganhos;
\item
  comportamento de market makers;
\item
  estabilidade e reflexividade;
\item
  capacidade de terceiros explorarem o mecanismo.
\end{itemize}

A região viável Θᵥ existe apenas se alguma configuração preservar
utilidade econômica para a companhia sem violar critérios mínimos de
integridade e funcionamento de mercado.

A tese não deve ser defendida contra o resultado de Θᵥ vazio.

\begin{center}\rule{0.5\linewidth}{0.5pt}\end{center}

\hypertarget{vi-evidencia-precedentes-e-validacao}{}

\hypertarget{vi-eviduxeancia-precedentes-e-validauxe7uxe3o-1}{%
\section{VI. Evidência, precedentes e
validação}\label{vi-eviduxeancia-precedentes-e-validauxe7uxe3o-1}}

\hypertarget{29-o-que-a-literatura-ja-permite-afirmar}{}

\hypertarget{29-o-que-a-literatura-juxe1-permite-afirmar}{%
\subsection{29. O que a literatura já permite
afirmar}\label{29-o-que-a-literatura-juxe1-permite-afirmar}}

A literatura não fornece sinal único para o efeito de recompras sobre
liquidez.

Brockman e Chung, usando mais de cinco mil recompras efetivamente
identificadas em Hong Kong, encontraram evidência de timing gerencial e
deterioração de algumas medidas de liquidez durante períodos de
recompra, incluindo spreads maiores e menor profundidade
(\protect\hyperlink{ref-7}{Brockman e Chung, 2001}).

Hillert, Maug e Obernberger, analisando 50.204 meses de recompra
realizada nos Estados Unidos entre 2004 e 2010, encontraram efeito
oposto em sua amostra: recompras podem melhorar liquidez, inclusive
porque a companhia atua como fornecedora de liquidez quando outros
investidores desejam vender (\protect\hyperlink{ref-8}{Hillert, Maug e
Obernberger, 2016}).

Ginglinger e Hamon encontram efeitos adversos sobre liquidez em dados
franceses e interpretam boa parte da atuação como contrarian trading e
suporte de preço (\protect\hyperlink{ref-9}{Ginglinger e Hamon, 2007}).

Portanto:

\begin{quote}
\textbf{o emissor pode simultaneamente ser fonte de seleção adversa e
fonte de liquidez. O sinal líquido é empírico e dependente do contexto.}
\end{quote}

Isso é importante porque a tese nunca deve assumir ``presença do emissor
reduz liquidez'' como premissa. Spread, depth, volume, impacto e
comportamento de market makers são outputs.

A literatura de microestrutura acrescenta outro ponto: trades carregam
informação e participantes ajustam preços e quotes em resposta ao risco
de negociar contra agentes informados (\protect\hyperlink{ref-4}{Glosten
e Milgrom, 1985}; \protect\hyperlink{ref-5}{Hasbrouck, 1991}). Isso
torna plausível tanto o benefício informacional quanto o risco de
amplificação do sinal identificado do emissor.

Uma análise exploratória dos Formulários de Referência da CVM enviados
para esta pesquisa também sugere que programas de recompra não são
meramente autorizações dormentes. Nos FREs de 2021 e 2022,
aproximadamente três quartos dos registros de programas com quantidade
prevista apresentavam alguma quantidade efetivamente adquirida; entre os
programas que executaram, a mediana de execução ficou próxima de dois
terços da quantidade prevista. As tabelas de movimentação de Treasury
também mostram aquisição e alienação efetivas por um conjunto material
de companhias. Esses dados são úteis para estabelecer
\textbf{familiaridade institucional e execução agregada}, mas não
permitem reconstruir de forma limpa e padronizada a frequência por
sessão. Programa relativamente comum não implica execução diária
relativamente comum.

Essa lacuna é relevante para estimar a saliência inicial do regime, mas
não invalida sua arquitetura conceitual. A frequência operacional sob o
protocolo será também parcialmente endógena a H, Q, T, à duração dos
programas e às decisões das próprias companhias.

\begin{center}\rule{0.5\linewidth}{0.5pt}\end{center}

\hypertarget{30-o-que-precedentes-regulatorios-sustentam-e-o-que-nao-sustentam}{}

\hypertarget{30-o-que-precedentes-regulatuxf3rios-sustentam--e-o-que-nuxe3o-sustentam}{%
\subsection{30. O que precedentes regulatórios sustentam --- e o que não
sustentam}\label{30-o-que-precedentes-regulatuxf3rios-sustentam--e-o-que-nuxe3o-sustentam}}

A proposta não surgiu de uma tentativa de copiar um regime existente.
Ainda assim, alguns componentes possuem paralelos institucionais úteis.

Nos Estados Unidos, a Rule 10b-18 da SEC estabelece condições de safe
harbour para recompras, incluindo restrição de volume. A SEC explica que
a limitação --- em geral 25\% do ADTV --- busca impedir que o emissor
domine o mercado de suas próprias ações e comprometa a percepção do
mercado como mecanismo independente de formação de preço
(\protect\hyperlink{ref-10}{SEC, 2002}).

Na União Europeia, o Regulamento Delegado 2016/1052 também estabelece
condições de preço e volume; o emissor não pode, para fins do safe
harbour, comprar em um dia mais de 25\% do volume médio diário no venue.
O regulamento exige reporte e divulgação das transações do programa até
o fim da sétima sessão de mercado após a execução
(\protect\hyperlink{ref-3}{European Commission, 2016}).

Esses regimes mostram que \textbf{capacidade quantitativa de
intervenção, timing de disclosure e perímetro instrumental são objetos
regulatórios reais}.

Eles não possuem nossa arquitetura completa.

Em particular:

\begin{itemize}
\tightlist
\item
  o limite de volume existente não é \(H_{\mathrm{gross}}\) no sentido
  adaptativo desta tese;
\item
  não há nosso \(H_{\mathrm{net}}\) direcional;
\item
  não há nossa função U → I → H futuro;
\item
  o disclosure europeu por prazo após execução é mais próximo de um
  trigger temporal do que de Q;
\item
  não há nosso Q como threshold de utilização que produz mudança de
  estado;
\item
  não há nosso H progressivo baseado em informação devolvida ao mercado.
\end{itemize}

Portanto, convergência de componentes não é identidade de mecanismo.

No Brasil, a própria CVM abriu a Consulta Pública SDM 04/2025 para
discutir ajustes na Resolução 77 com foco em negociação de ações de
própria emissão, apoiada em análise de impacto regulatório sobre
recompra e liquidez e em comparações internacionais
(\protect\hyperlink{ref-11}{CVM, 2025}). Isso não valida a proposta, mas
confirma a atualidade institucional do problema.

\begin{center}\rule{0.5\linewidth}{0.5pt}\end{center}

\hypertarget{31-casos-administrativos-como-testes-de-realidade}{}

\hypertarget{31-casos-administrativos-como-testes-de-realidade}{%
\subsection{31. Casos administrativos como testes de
realidade}\label{31-casos-administrativos-como-testes-de-realidade}}

Casos reais são úteis não para provar a tese, mas para impedir que ela
seja construída sobre uma imagem excessivamente limpa do mercado. O
exercício é contrafactual: primeiro se apresenta o que foi publicamente
discutido no processo; depois se pergunta quais informações, limites e
registros existiriam se o protocolo estivesse disponível. O resultado
jurídico histórico não é reescrito, e a aplicação hipotética do
protocolo não presume culpa nem inocência.

\textbf{Kepler Weber.} A área técnica da CVM propôs responsabilização da
companhia e de seu diretor de relações com investidores pela aquisição,
em nome da companhia, de 112 mil ações de própria emissão entre 13 e 27
de julho de 2022, dentro do período de vedação anterior à divulgação do
ITR. O processo foi encerrado por Termo de Compromisso, sem julgamento
de mérito (\protect\hyperlink{ref-12}{CVM, 2023a}). A distribuição das
aquisições por múltiplos pregões mostra por que a unidade temporal
economicamente relevante não precisa coincidir com um único dia.

No universo do protocolo, a primeira aquisição abriria W; H limitaria a
escala disponível naquela janela; Q ou T determinariam o momento de
devolução informacional; e eventual continuidade ocorreria apenas em
janelas posteriores. O mercado observaria que a própria Kepler Weber
estava comprando, sem receber automaticamente a informação privada que
motivava a decisão. O caso deixa de ser apenas um problema binário de
permissão ou proibição e passa a ser também um teste de capacidade,
duração, disclosure e reconstrução da trajetória.

\textbf{JBS, FB Participações e J\&F.} O processo examinou negócios
realizados pela JBS S.A. e por sua então controladora, FB Participações
S.A., com ações JBSS3. Os acusados eram Joesley Batista, Wesley Batista
e J\&F Investimentos S.A., sucessora da FB Participações; a JBS era
parte central dos fatos, mas não figurava entre os acusados desse
processo. A decisão condenou a J\&F, nessa qualidade de sucessora,
porque a FB vendeu ações da JBS durante período vedado pelo programa de
recompra da companhia. Joesley, Wesley e J\&F foram absolvidos das
imputações de insider trading e manipulação relacionadas às operações
examinadas (\protect\hyperlink{ref-13}{CVM, 2023b}). A distinção é
precisamente o que torna o caso útil: negociação, informação, intenção e
manipulação não podem ser inferidas umas das outras de maneira
automática.

Sob o protocolo, a identificação do emissor, o consumo de H, os alertas
de Q, o fechamento de W e a trilha de execuções tornariam observável uma
parcela maior da conduta econômica. Isso não resolveria por fórmula a
existência de informação privilegiada ou intenção manipulativa, nem
afastaria proibições jurídicas externas ao regime. Permitiria, contudo,
separar com mais clareza o que foi solicitado, executado, divulgado e
economicamente capturado.

\textbf{JSL e Haitong.} Em processo relacionado a negócios com JSLG3
realizados durante programa de recompra, a área técnica da CVM formulou
acusações de manipulação contra a companhia e participantes ligados à
execução, além de prática não equitativa atribuída a participantes da
Haitong que teriam se antecipado à atuação esperada da JSL. O processo
foi objeto de Termos de Compromisso; esse desfecho não equivale a
julgamento de mérito ou admissão das acusações
(\protect\hyperlink{ref-14}{CVM, 2018}).

O caso mostra que o fluxo esperado do emissor pode tornar-se objeto
econômico para terceiros. No protocolo, a atuação identificada não
eliminaria vazamento, coordenação ou antecipação indevida; criaria,
porém, uma referência pública comum e uma trilha que separa atividade
Gross, posição Net, contrapartes, timing e reação. Também mostra por que
\textbf{Net = 0 não elimina relevância de Gross}: compras e vendas podem
se compensar na posição final e ainda assim produzir intervenção, sinal
e oportunidade de exploração.

O comparativo contrafactual pode ser resumido assim:

\begin{longtable}[]{@{}
  >{\raggedright\arraybackslash}p{(\columnwidth - 4\tabcolsep) * \real{0.1600}}
  >{\raggedright\arraybackslash}p{(\columnwidth - 4\tabcolsep) * \real{0.3600}}
  >{\raggedright\arraybackslash}p{(\columnwidth - 4\tabcolsep) * \real{0.4800}}@{}}
\toprule\noalign{}
\begin{minipage}[b]{\linewidth}\raggedright
Caso
\end{minipage} & \begin{minipage}[b]{\linewidth}\raggedright
Problema que tensiona a arquitetura
\end{minipage} & \begin{minipage}[b]{\linewidth}\raggedright
O que o protocolo tornaria observável ou limitado
\end{minipage} \\
\midrule\noalign{}
\endhead
\bottomrule\noalign{}
\endlastfoot
Kepler Weber & recompras distribuídas por múltiplos pregões & abertura e
encerramento de W, consumo de H e disclosures sucessivos \\
JBS e FB/J\&F & dificuldade de separar negociação, informação e intenção
ex post & identidade, utilização, trajetória executada e momento dos
disclosures, sem automatizar o julgamento jurídico \\
JSL/Haitong & manipulação alegada e antecipação do fluxo do emissor por
terceiros & Gross, Net, timing, reação e trilha auditável para
investigar exploração e coordenação \\
\end{longtable}

Os casos reforçam duas perguntas do laboratório:

\begin{quote}
terceiros conseguem capturar sistematicamente valor antecipando a
companhia?
\end{quote}

\begin{quote}
a companhia consegue usar a reação de terceiros ao próprio sinal como
fonte de retorno independente da informação fundamental?
\end{quote}

\begin{center}\rule{0.5\linewidth}{0.5pt}\end{center}

\hypertarget{32-estrategia-de-validacao-e-laboratorio-multiagente}{}

\hypertarget{32-estratuxe9gia-de-validauxe7uxe3o-e-laboratuxf3rio-multiagente}{%
\subsection{32. Estratégia de validação e laboratório
multiagente}\label{32-estratuxe9gia-de-validauxe7uxe3o-e-laboratuxf3rio-multiagente}}

A validação deve ocorrer em camadas.

\hypertarget{32-1-testes-de-invariantes}{}

\hypertarget{321-testes-de-invariantes}{%
\subsubsection{32.1. Testes de
invariantes}\label{321-testes-de-invariantes}}

Primeiro, verificar propriedades mecânicas: Treasury nunca negativa,
SELL limitado ao estoque, Gross não apagado por reversão, Q não
excedendo H, execução aceita nunca ultrapassando os limites
simultaneamente aplicáveis, residual excedente sendo rejeitado e
registrado, alertas intraday permanecendo auditáveis, disclosure
quantitativo sendo decidido pela condição consolidada no corte, Q ou T
encerrando W, capacidade não sendo reciclada dentro da janela encerrada
e a próxima W tornando-se elegível apenas a partir do pregão seguinte.

\hypertarget{32-2-testes-de-trajetoria}{}

\hypertarget{322-testes-de-trajetuxf3ria}{%
\subsubsection{32.2. Testes de
trajetória}\label{322-testes-de-trajetuxf3ria}}

Depois, estudar o mecanismo sem agentes sofisticados:

\begin{itemize}
\tightlist
\item
  BUY persistente;
\item
  SELL persistente;
\item
  alternância;
\item
  churn;
\item
  inatividade;
\item
  decay;
\item
  saturação;
\item
  proximidade dos extremos de Treasury;
\item
  Q baixo versus Q alto;
\item
  H₀ próximo versus distante de \(H_{\max}\);
\item
  crescimento linear versus acelerado de H;
\item
  reversão com e sem degradação de capacidade;
\item
  plano público com e sem bônus inicial;
\item
  regimes de baixa e alta familiaridade operacional.
\end{itemize}

É nessa fase que um exemplo numérico completo deve ser construído. Ele
não será apresentado como calibração realista, mas como \textbf{exemplo
pedagógico executável} para transformar a sopa de símbolos em uma
trajetória concreta: H inicial, execução, U, alerta Q, consolidação no
fechamento, decisão de disclosure, atualização de I, decay e novo H.

\hypertarget{32-3-implementacao-de-referencia}{}

\hypertarget{323-implementauxe7uxe3o-de-referuxeancia-issuer-protocol-lab}{%
\subsubsection{32.3. Implementação de referência: Issuer Protocol
Lab}\label{323-implementauxe7uxe3o-de-referuxeancia-issuer-protocol-lab}}

O \textbf{Issuer Protocol Lab} implementa a primeira camada
determinística da estratégia de validação. Sua função é tornar a máquina
regulatória executável, comparável e auditável antes da introdução de
microestrutura ou agentes adaptativos. Ele não demonstra que a
arquitetura é desejável, não estima uma calibração ótima e não substitui
validação empírica.

O Lab organiza o experimento em cinco tipos de objeto:

\begin{longtable}[]{@{}
  >{\raggedright\arraybackslash}p{(\columnwidth - 2\tabcolsep) * \real{0.1800}}
  >{\raggedright\arraybackslash}p{(\columnwidth - 2\tabcolsep) * \real{0.8200}}@{}}
\toprule\noalign{}
\begin{minipage}[b]{\linewidth}\raggedright
Objeto
\end{minipage} & \begin{minipage}[b]{\linewidth}\raggedright
Função
\end{minipage} \\
\midrule\noalign{}
\endhead
\bottomrule\noalign{}
\endlastfoot
dados & snapshot observado da companhia, com identificação e
proveniência \\
variáveis & transformações declaradas dos dados primitivos ou de outras
variáveis \\
cenário & trajetória exógena de ordens de uma companhia \\
especificação & regras de memória, capacidade e intrawindow \\
avaliação & resultado elementar de uma combinação especificação ×
cenário; comparativos reúnem várias avaliações mantendo um dos domínios
fixo \\
\end{longtable}

Cada cenário registra companhia, estratégia e sequência de ordens
solicitadas. Cada especificação reúne fórmulas e parâmetros
substituíveis para sinal, atualização da memória, decay, H₀,
\(H_{\max}\), saturação e Q/H. A separação é deliberada: cenários podem
ser reutilizados entre especificações, e especificações podem ser
avaliadas contra diferentes cenários sem incorporar a estratégia de
ordens à regra regulatória.

Uma avaliação determinística segue o ciclo:

\begin{quote}
snapshot e mapeamentos semânticos → primeira execução e abertura de W →
cálculo e trava dos H da janela → submissão das ordens → execução ou
rejeição parcial → alertas intraday de Q → consolidação no fechamento →
encerramento por Q ou T e disclosure → atualização do estado
\end{quote}

Os H permanecem travados durante a janela avaliada. Cada ordem é
confrontada simultaneamente com \(H_{\mathrm{gross}}\),
\(H_{\mathrm{net}+}\), \(H_{\mathrm{net}-}\) e as restrições físicas
aplicáveis. Quando a quantidade solicitada não é integralmente
admissível, o engine executa a parcela possível e registra o residual
rejeitado e sua causa. A trajetória regulatória contém apenas
quantidades executadas; intenções rejeitadas permanecem disponíveis para
auditoria, mas não alteram o estado.

O comparativo implementa dois contrafactuais controlados. Em \textbf{uma
especificação contra vários cenários}, cada trajetória é normalizada
pelo H resolvido para a própria companhia, permitindo comparar
utilização relativa; os valores absolutos permanecem como metadados. Em
\textbf{um cenário contra várias especificações}, o eixo utiliza
quantidades absolutas de ações, porque cada especificação produz seus
próprios H e Q sobre a mesma trajetória solicitada. A simetria
metodológica consiste em manter um domínio fixo e variar o outro, não em
impor a mesma normalização aos dois casos.

Cada resultado preserva a configuração utilizada, a trajetória
solicitada, a trajetória executada, os alertas, a decisão de disclosure,
as parcelas rejeitadas e os invariantes avaliados. Especificações
concluídas recebem identidade reproduzível; alterações posteriores devem
permanecer distinguíveis de execuções anteriores. Essa rastreabilidade
permite reconstruir por que uma ordem foi executada, limitada ou
rejeitada e quais regras produziram cada capacidade.

Os dados, cenários, fórmulas, parâmetros e especificações distribuídos
previamente com o Lab são \textbf{fixtures de demonstração e teste
visual}. Eles existem para exercitar caminhos do sistema, tornar
diferenças observáveis e revelar inconsistências mecânicas. Não
constituem estimativas empíricas, calibração plausível, recomendação
regulatória ou evidência de que os valores escolhidos pertençam a uma
região viável. Os dados atualmente disponíveis são insuficientes para
sustentar qualquer dessas interpretações.

O estágio implementado deve, portanto, ser entendido como validação de
consistência interna. Ele permite testar invariantes, semântica
temporal, composição de fórmulas, limites físicos, reprodutibilidade e
trajetórias extremas. Não permite concluir sobre spread, depth, impacto
de preço, descoberta descentralizada, comportamento estratégico, welfare
ou equilíbrio de mercado. Essas perguntas pertencem às camadas
seguintes.

\hypertarget{32-4-variable-h-lab}{}

\hypertarget{324-variableh-lab}{%
\subsubsection{32.4. Variable/H Lab}\label{324-variableh-lab}}

O laboratório de variáveis permite comparar especificações de H mantendo
o restante do mecanismo constante. Sua função é descobrir quais
componentes acrescentam valor e quais apenas aumentam complexidade.

\hypertarget{32-5-microestrutura}{}

\hypertarget{325-microestrutura}{%
\subsubsection{32.5. Microestrutura}\label{325-microestrutura}}

Em seguida entram book, spread, depth, impacto, market makers e
diferentes condições de liquidez.

\hypertarget{32-6-agentes-adaptativos}{}

\hypertarget{326-agentes-adaptativos}{%
\subsubsection{32.6. Agentes
adaptativos}\label{326-agentes-adaptativos}}

Os agentes adaptativos ainda \textbf{não integram a implementação de
referência}. Faltam dados públicos suficientemente granulares para
estimar suas heurísticas e falta um baseline regulatório observado que
permita distinguir aprendizagem plausível de comportamento
arbitrariamente programado. Essa ausência deve ser tratada como limite
atual da validação, não escondida como funcionalidade futura já
resolvida.

Ainda assim, um laboratório multiagente é um instrumento possível e
importante para testar a regulação antes de expor o mercado real ao
mecanismo. Em vez de ``lançar a bomba'' e aprender apenas depois da
adoção, o regulador pode construir companhias, market makers,
investidores informados e menos informados, arbitradores, traders
direcionais, terceiros sofisticados e um regulador fictícios,
submetendo-os a diferentes especificações e condições de mercado.

Cada agente deve possuir heurísticas explícitas e capacidade de aprender
com o regime. O objetivo não é apenas encontrar uma estratégia
vencedora, mas entender \textbf{por que ela surgiu, qual parâmetro
permitiu sua exploração e como o equilíbrio mudou depois que outros
agentes aprenderam}.

Uma camada adversarial deve aproximar o caso limite de agentes altamente
racionais: a companhia procura a política ótima dadas as regras e a
reação esperada do mercado; participantes procuram a resposta ótima
dadas as ações da companhia e as restrições que o protocolo lhe impõe. A
dinâmica é iterativa: companhia adapta → mercado adapta → companhia
readapta. O objetivo não é assumir racionalidade perfeita como descrição
empírica, mas utilizá-la como stress test contra exploits estruturais.

O laboratório deve incluir cold start e adoção gradual. Um regime
inicialmente raro pode produzir sinal muito mais saliente do que um
regime já familiar. A validação deve observar se aprendizagem reduz
overreaction, se market makers retiram liquidez antes da estabilização,
se reputações heterogêneas geram equilíbrios distintos e se a companhia
consegue realizar valor após repricing sem transformar o mecanismo em
estratégia mecânica de fabricação de sinal.

O regulador também é experimentador. Ele compara configurações, não
apenas pontos isolados, observa fronteiras de trade-off e procura
configurações em que nenhuma estratégia exploratória consiga escala
relevante sem pagar em capacidade, risco, observabilidade ou
opcionalidade futura.

A implementação faz parte da metodologia da tese porque força decisões
semânticas que o texto pode esconder. Se implementar o protocolo exigir
exceções ad hoc que não derivam da arquitetura, isso é evidência de que
a formalização ainda possui lacunas.

O laboratório integra a estratégia de validação da proposta: seus
resultados relevantes devem ser incorporados à avaliação empírica do
mecanismo, em vez de tratados como etapa externa à tese.

\begin{center}\rule{0.5\linewidth}{0.5pt}\end{center}

\hypertarget{vii-conclusao}{}

\hypertarget{vii-conclusuxe3o-1}{%
\section{VII. Conclusão}\label{vii-conclusuxe3o-1}}

\hypertarget{33-o-que-esta-tese-estabelece}{}

\hypertarget{33-o-que-esta-tese-estabelece}{%
\subsection{33. O que esta tese
estabelece}\label{33-o-que-esta-tese-estabelece}}

Esta tese não estabelece que o mecanismo é superior ao regime vigente.

Sua hipótese normativa é justamente investigar se a companhia pode
receber liberdade para negociar mesmo quando possui informação privada,
sem produzir lesão sistemática aos demais participantes. Essa liberdade
não significa ausência de regras: significa dispensar uma proibição
geral ou uma autorização individual, analisada caso a caso, quando a
operação respeita W, H, Q, T, Treasury, disclosure e os demais
invariantes do regime. A tese não presume que essa compatibilidade
exista; ela constrói o teste capaz de aceitá-la ou rejeitá-la.

Não estabelece que disclosure elimina manipulação.

Não estabelece que lucro da Treasury seja sinônimo de welfare.

Não estabelece que maior observabilidade externa justifique
automaticamente H maior.

Não estabelece que o modelo calendarizado seja superior ao rolling ou
vice-versa.

Não estabelece valores ótimos para H, Q, T, decay ou qualquer outro
parâmetro.

Ela estabelece algo mais limitado:

\begin{quote}
\textbf{é possível decompor o problema da negociação própria do emissor
sob assimetria informacional em uma arquitetura explícita de capacidade,
posição, intervenção, utilização, disclosure, memória e aprendizagem
suficientemente precisa para ser testada e rejeitada.}
\end{quote}

A decomposição central é:

\begin{quote}
Gross → intensidade de intervenção
\end{quote}

\begin{quote}
Net → alteração de posição
\end{quote}

\begin{quote}
W → janela regulatória comum a H, Q e T
\end{quote}

\begin{quote}
H → capacidade máxima
\end{quote}

\begin{quote}
Q → limiar de alerta e avaliação quantitativa no corte
\end{quote}

\begin{quote}
U → intensidade relativa da utilização
\end{quote}

\begin{quote}
\(I_{\mathrm{issuer}}\) → memória informacional revelada pela ação
observável do emissor
\end{quote}

\begin{quote}
\(I_{\mathrm{market}}\) → informação economicamente relevante já
observável externamente
\end{quote}

\begin{quote}
decay → envelhecimento da informação
\end{quote}

\begin{quote}
Treasury → restrição física da posição
\end{quote}

\begin{quote}
reputação → interpretação endógena do mercado, não parâmetro oficial
\end{quote}

\begin{quote}
segurança jurídica → previsibilidade ex ante para negociação de boa-fé
dentro do regime, sem imunidade a manipulação
\end{quote}

A partir dessa ontologia, diferentes especificações podem ser comparadas
sem mudar a pergunta econômica.

\begin{center}\rule{0.5\linewidth}{0.5pt}\end{center}

\hypertarget{34-conclusao-conceitual}{}

\hypertarget{34-conclusuxe3o-conceitual}{%
\subsection{34. Conclusão conceitual}\label{34-conclusuxe3o-conceitual}}

A proposta parte de uma inversão simples.

A companhia não precisa deixar de conhecer melhor a si mesma para que o
mercado seja protegido. Talvez seja possível permitir que ela utilize
parte dessa vantagem, desde que a capacidade seja limitada e sua
utilização produza informação observável em ritmo suficiente para
impedir exclusividade indefinida.

\begin{quote}
\textbf{A vantagem não é o problema. O problema é a forma como ela é
exercida.}
\end{quote}

A companhia pode ganhar se suas decisões forem boas e pode perder se
forem ruins. O ganho pretendido não é permitir que ela capture sozinha
todo o movimento, mas preservar a possibilidade de obter \textbf{parte}
da valorização que ajudou a antecipar enquanto o mercado recebe o sinal,
reage e participa do restante da descoberta de preço. Decisões
bem-sucedidas podem preservar mais recursos corporativos para decisões
futuras; decisões ruins consomem patrimônio e deterioram credibilidade.

Essa consequência só é aceitável com a segunda metade da arquitetura.

A companhia não recebe o direito de transformar informação privada em
vantagem ilimitada; não recebe o direito de apagar intervenção por meio
de reversão; não recebe derivativos para alavancar o próprio sinal; e
não recebe garantia de que o mercado permanecerá passivo enquanto ela
opera. Ao tentar ampliar a captura, consome H, aproxima-se de Q,
antecipa o encerramento de W e revela mais da própria decisão. A
arquitetura procura transformar escala em informação antes que o emissor
consiga excluir o mercado do ajuste inteiro.

Ao contrário, sua atuação passa a produzir um histórico que outros
participantes podem observar, interpretar, antecipar e eventualmente
explorar.

O terceiro ganho pretendido é jurídico: uma companhia que negocia de
boa-fé dentro do regime deve poder conhecer ex ante as condições
operacionais que tornam sua atuação admissível, sem que o simples fato
de conhecer melhor a própria realidade econômica transforme toda decisão
de timing em uma zona de incerteza. Essa previsibilidade não protege
fraude ou manipulação; ela delimita o espaço legítimo de atuação.

Isso cria uma tensão deliberada:

\begin{quote}
\textbf{a companhia pode ser a primeira a saber e a primeira a agir ---
mas não pode ser a única a participar da descoberta de preço que sua
informação produz.}
\end{quote}

O desenho também assume explicitamente que nenhum conjunto finito de
regras elimina toda margem estratégica. Sua premissa é mais exigente e
mais realista: \textbf{nenhuma estratégia exploratória deve conseguir
escala relevante sem consumir capacidade, assumir risco econômico,
deixar rastro observável ou destruir a própria capacidade futura.} O
mercado participa dessa disciplina ao aprender, antecipar, contrariar e
precificar o comportamento do emissor.

A maior ameaça à proposta talvez seja justamente o sucesso desse
mecanismo informacional. Se a ação identificada do emissor se tornar um
sinal excessivamente dominante, o protocolo pode deteriorar liquidez,
criar reflexividade, reduzir produção independente de informação ou
permitir que reputação substitua fundamentos no curto prazo.

Por isso a pergunta final não é ``como calibrar H para funcionar?''.

É:

\begin{quote}
\textbf{Existe alguma região de desenho em que a companhia preserve
capacidade economicamente relevante e juridicamente previsível de agir,
o mercado preserve capacidade economicamente relevante de descobrir
preço por conta própria, e a informação revelada pela atuação do emissor
ajude mais do que desorganize esse processo?}
\end{quote}

Se a resposta experimental for não, a arquitetura deve ser rejeitada.

Se for sim, então a negociação própria do emissor pode deixar de ser
tratada apenas como fonte de conflito e passar também a ser investigada
como instrumento de alocação de capital e produção de informação --- sem
exigir que a assimetria informacional desapareça antes que qualquer
decisão econômica possa ocorrer.

Essa é a hipótese que o laboratório precisa tentar destruir.

\begin{center}\rule{0.5\linewidth}{0.5pt}\end{center}

\hypertarget{35-referencias}{}

\hypertarget{35-referuxeancias}{%
\subsection{35. Referências}\label{35-referuxeancias}}

\hypertarget{ref-1}{} \textbf{Comissão de Valores Mobiliários (2022).}
Resolução CVM nº 77, de 29 de março de 2022 --- negociação de ações e
aquisição de debêntures de própria emissão.\\
\href{https://conteudo.cvm.gov.br/legislacao/resolucoes/resol077.html}{Acessar
fonte}

\hypertarget{ref-2}{} \textbf{Comissão de Valores Mobiliários (2021).}
Resolução CVM nº 44, de 23 de agosto de 2021 --- ato ou fato relevante e
negociação na pendência de informação relevante não divulgada.\\
\href{https://conteudo.cvm.gov.br/legislacao/resolucoes/resol044.html}{Acessar
fonte}

\hypertarget{ref-3}{} \textbf{European Commission (2016).} Commission
Delegated Regulation (EU) 2016/1052, especialmente arts. 2 e 3 ---
disclosure, preço e limite de 25\% do volume médio diário para o safe
harbour de recompra.\\
\href{https://eur-lex.europa.eu/eli/reg_del/2016/1052}{Acessar fonte}

\hypertarget{ref-4}{} \textbf{Glosten, L. R.; Milgrom, P. R. (1985).}
\emph{Bid, ask and transaction prices in a specialist market with
heterogeneously informed traders}. Journal of Financial Economics,
14(1), 71--100.\\
\href{https://doi.org/10.1016/0304-405X(85)90044-3}{Acessar fonte}

\hypertarget{ref-5}{} \textbf{Hasbrouck, J. (1991).} \emph{Measuring the
Information Content of Stock Trades}. The Journal of Finance, 46(1),
179--207.\\
\href{https://doi.org/10.1111/j.1540-6261.1991.tb03749.x}{Acessar fonte}

\hypertarget{ref-6}{} \textbf{Grossman, S. J.; Stiglitz, J. E. (1980).}
\emph{On the Impossibility of Informationally Efficient Markets}.
American Economic Review, 70(3), 393--408.\\
\href{https://www.jstor.org/stable/1805228}{Acessar fonte}

\hypertarget{ref-7}{} \textbf{Brockman, P.; Chung, D. Y. (2001).}
\emph{Managerial timing and corporate liquidity: evidence from actual
share repurchases}. Journal of Financial Economics, 61(3), 417--448.

\hypertarget{ref-8}{} \textbf{Hillert, A.; Maug, E.; Obernberger, S.
(2016).} \emph{Stock Repurchases and Liquidity}. Journal of Financial
Economics, 119(1), 186--209.

\hypertarget{ref-9}{} \textbf{Ginglinger, E.; Hamon, J. (2007).}
\emph{Actual share repurchases, timing and liquidity}. Journal of
Banking \& Finance, 31(3), 915--938.

\hypertarget{ref-10}{} \textbf{U.S. Securities and Exchange Commission
(2002).} Rule 10b-18 and Purchases of Certain Equity Securities by the
Issuer and Others --- discussão da condição de volume e do limite de
25\% do ADTV.\\
\href{https://www.sec.gov/rules-regulations/2002/12/rule-10b-18-purchases-certain-equity-securities-issuer-others}{Acessar
fonte}

\hypertarget{ref-11}{} \textbf{Comissão de Valores Mobiliários (2025).}
Consulta Pública SDM 04/2025 --- ajustes na Resolução CVM 77 com foco em
negociação de ações de própria emissão.\\
\href{https://conteudo.cvm.gov.br/cvm_institucional/audiencias_publicas/ap_sdm/2025/sdm0425.html}{Acessar
fonte}

\hypertarget{ref-12}{} \textbf{Comissão de Valores Mobiliários (2023a).}
Termo de Compromisso no PAS CVM 19957.014436/2022-29 --- aquisição de
ações de própria emissão pela Kepler Weber S.A. no período anterior à
divulgação do ITR.\\
\href{https://www.gov.br/cvm/pt-br/assuntos/noticias/2023/cvm-aceita-proposta-conjunta-de-mais-de-r-900-mil-apresentada-por-kepler-weber-s-a-e-seu-diretor-de-relacoes-com-investidores}{Acessar
fonte}

\hypertarget{ref-13}{} \textbf{Comissão de Valores Mobiliários (2023b).}
Julgamento de processos envolvendo empresas do grupo JBS e seus
executivos --- decisões sobre período vedado, uso de informação
privilegiada e manipulação.\\
\href{https://www.gov.br/cvm/pt-br/assuntos/noticias/2023/cvm-conclui-julgamento-de-processos-envolvendo-empresas-do-grupo-jbs-e-seus-executivos}{Acessar
fonte}

\hypertarget{ref-14}{} \textbf{Comissão de Valores Mobiliários (2018).}
Apreciação de propostas de Termo de Compromisso no PAS
19957.010277/2017-26 --- operações com JSLG3 durante programa de
recompra e antecipação atribuída a participantes da Haitong.\\
\href{https://conteudo.cvm.gov.br/decisoes/2018/20180814_R1/20180814_D1023.html}{Acessar
fonte}

\begin{center}\rule{0.5\linewidth}{0.5pt}\end{center}

\hypertarget{nota-epistemoluxf3gica}{%
\subsubsection{Nota epistemológica}\label{nota-epistemoluxf3gica}}

As referências acima cumprem funções diferentes. Literatura de
microestrutura sustenta a plausibilidade de trades carregarem informação
e de seleção adversa afetar formação de preços. Estudos de recompra
mostram que o sinal de liquidez não é único. Regras dos Estados Unidos,
União Europeia e Brasil demonstram que volume, timing, disclosure e
negociação própria são objetos regulatórios reais. Nenhuma dessas fontes
demonstra que a arquitetura proposta seja superior ao regime vigente ou
que exista Θᵥ não vazio.

A validade institucional da proposta depende, portanto, de confrontar a
arquitetura com implementação, dados e comportamento estratégico.

\end{document}
